\pdfoutput=1
% ---------------------------------------------------------------------------
% Does coupling divergence tell a god unit from a complex algorithm?
% A pre-registered test of weighting complexity concentration by the data a
% unit hands on.
%
% Built for arXiv (primary class cs.SE). No .bib: the bibliography is a
% `thebibliography` environment at the end.
%
% Every number in this file is a macro from generated/numbers.tex and every
% table is a fragment under generated/, both written by paper/scripts/tables.py
% from the harness's confirmatory analysis. No number this study computes is
% typed by hand; the only typed numbers are literature constants and the
% plan's own constants. Code, commits, files and commands are named in the
% ancillary README, not here.
%
% Build:   pdflatex main && pdflatex main
% ---------------------------------------------------------------------------
\documentclass[11pt,a4paper]{article}

\usepackage[T1]{fontenc}
\usepackage[utf8]{inputenc}
\usepackage{lmodern}
\usepackage[margin=2.4cm]{geometry}
\usepackage{graphicx}
\usepackage{booktabs}
\usepackage{multirow}
\usepackage{amsmath}
\usepackage{microtype}
\usepackage[hidelinks]{hyperref}
\usepackage{caption}
\captionsetup{font=small,labelfont=bf}

% generated by paper/scripts/tables.py; do not edit
\newcommand{\analysisMadeUtc}{2026-10-26 14:55:06 UTC}
\newcommand{\analysisVOneMadeUtc}{2026-10-26 08:04:12 UTC}
\newcommand{\analysisKind}{confirmatory}
\newcommand{\analysisShaShort}{9aa901ddc12b}
\newcommand{\analysisVOneShaShort}{d64d4c8d547f}
\newcommand{\decisionsMadeUtc}{2026-10-26 13:53:46 UTC}
\newcommand{\bootB}{2{,}000}
\newcommand{\bootSeed}{12345}
\newcommand{\mlcqRepos}{522}
\newcommand{\mlcqSamples}{4{,}764}
\newcommand{\mlcqBlobReviewed}{2{,}334}
\newcommand{\mlcqBlob}{254}
\newcommand{\mlcqLongReviewed}{2{,}430}
\newcommand{\mlcqLong}{276}
\newcommand{\mlcqMovedMeasured}{23}
\newcommand{\mlcqMovedConfirm}{12}
\newcommand{\mlcqMovedWithBlob}{12}
\newcommand{\mlcqMeasured}{520}
\newcommand{\mlcqConfirmMeasured}{261}
\newcommand{\mlcqConfirmUsable}{258}
\newcommand{\mlcqConfirmErrors}{2}
\newcommand{\mlcqExploreMeasured}{261}
\newcommand{\seedThreeMeasured}{500}
\newcommand{\seedThreeUsable}{492}
\newcommand{\seedThreeDrawUsable}{580}
\newcommand{\seedThreeExcluded}{62}
\newcommand{\fontanaSystems}{73}
\newcommand{\fontanaInstances}{420}
\newcommand{\fontanaGodInstances}{140}
\newcommand{\fontanaFilesReviewed}{384}
\newcommand{\fontanaFilesGod}{139}
\newcommand{\fontanaFilesReviewedMeasured}{380}
\newcommand{\fontanaFilesGodMeasured}{139}
\newcommand{\validationRepos}{150}
\newcommand{\joinMlcqLizardUnits}{1{,}949{,}548}
\newcommand{\joinMlcqShare}{99.4\%}
\newcommand{\joinMlcqMeasurerExtra}{40{,}390}
\newcommand{\joinMlcqMeasurerExtraShare}{2.0\%}
\newcommand{\joinMlcqClassUnits}{179{,}585}
\newcommand{\joinMlcqClassCoupledShare}{99.9\%}
\newcommand{\callsMlcqInternal}{42.3\%}
\newcommand{\callsMlcqJdk}{25.8\%}
\newcommand{\callsMlcqUnresolved}{24.0\%}
\newcommand{\callsMlcqShouldHave}{7.9\%}
\newcommand{\callsMlcqBudget}{0.0\%}
\newcommand{\roleMlcqClassCatchmentUnits}{50.2\%}
\newcommand{\roleMlcqClassCatchmentCc}{45.5\%}
\newcommand{\roleMlcqClassPipeUnits}{6.6\%}
\newcommand{\roleMlcqClassPipeCc}{1.4\%}
\newcommand{\roleMlcqClassJunctionUnits}{43.2\%}
\newcommand{\roleMlcqClassJunctionCc}{53.1\%}
\newcommand{\roleMlcqMethodCatchmentUnits}{56.6\%}
\newcommand{\roleMlcqMethodCatchmentCc}{40.9\%}
\newcommand{\roleMlcqMethodPipeUnits}{12.4\%}
\newcommand{\roleMlcqMethodPipeCc}{9.1\%}
\newcommand{\roleMlcqMethodJunctionUnits}{31.0\%}
\newcommand{\roleMlcqMethodJunctionCc}{50.0\%}
\newcommand{\roleMlcqMethodDataCatchmentUnits}{35.2\%}
\newcommand{\roleMlcqMethodDataPipeUnits}{40.0\%}
\newcommand{\throwNewMlcqClassShare}{2.9\%}
\newcommand{\throwNewMlcqMethodShare}{2.1\%}
\newcommand{\diMlcqRepos}{113}
\newcommand{\prunedMlcqFiles}{1{,}271}
\newcommand{\filesMlcq}{486{,}712}
\newcommand{\filesMlcqFailed}{1{,}251}
\newcommand{\filesMlcqFailedShare}{0.26\%}
\newcommand{\scaleMlcqClassRepos}{259}
\newcommand{\scaleMlcqXMad}{257}
\newcommand{\scaleMlcqDMad}{94}
\newcommand{\scaleMlcqDMeanAd}{163}
\newcommand{\scaleMlcqDConstant}{2}
\newcommand{\zeroMlcqDClass}{57.3\%}
\newcommand{\scaleMlcqMethodRepos}{259}
\newcommand{\scaleMlcqXmMeanAd}{244}
\newcommand{\scaleMlcqDmMeanAd}{259}
\newcommand{\zeroMlcqXm}{63.5\%}
\newcommand{\zeroMlcqDmMethod}{80.5\%}
\newcommand{\joinSeedThreeLizardUnits}{2{,}993{,}414}
\newcommand{\joinSeedThreeShare}{98.3\%}
\newcommand{\joinSeedThreeMeasurerExtra}{89{,}757}
\newcommand{\joinSeedThreeMeasurerExtraShare}{3.0\%}
\newcommand{\joinSeedThreeClassUnits}{295{,}031}
\newcommand{\joinSeedThreeClassCoupledShare}{98.5\%}
\newcommand{\callsSeedThreeInternal}{37.7\%}
\newcommand{\callsSeedThreeJdk}{25.1\%}
\newcommand{\callsSeedThreeUnresolved}{27.3\%}
\newcommand{\callsSeedThreeShouldHave}{9.7\%}
\newcommand{\callsSeedThreeBudget}{0.0\%}
\newcommand{\roleSeedThreeClassCatchmentUnits}{50.7\%}
\newcommand{\roleSeedThreeClassCatchmentCc}{42.8\%}
\newcommand{\roleSeedThreeClassPipeUnits}{6.4\%}
\newcommand{\roleSeedThreeClassPipeCc}{1.5\%}
\newcommand{\roleSeedThreeClassJunctionUnits}{42.9\%}
\newcommand{\roleSeedThreeClassJunctionCc}{55.7\%}
\newcommand{\roleSeedThreeMethodCatchmentUnits}{57.0\%}
\newcommand{\roleSeedThreeMethodCatchmentCc}{41.5\%}
\newcommand{\roleSeedThreeMethodPipeUnits}{11.9\%}
\newcommand{\roleSeedThreeMethodPipeCc}{9.1\%}
\newcommand{\roleSeedThreeMethodJunctionUnits}{31.1\%}
\newcommand{\roleSeedThreeMethodJunctionCc}{49.4\%}
\newcommand{\roleSeedThreeMethodDataCatchmentUnits}{36.4\%}
\newcommand{\roleSeedThreeMethodDataPipeUnits}{38.6\%}
\newcommand{\throwNewSeedThreeClassShare}{2.2\%}
\newcommand{\throwNewSeedThreeMethodShare}{1.6\%}
\newcommand{\diSeedThreeRepos}{163}
\newcommand{\prunedSeedThreeFiles}{26{,}123}
\newcommand{\filesSeedThree}{359{,}093}
\newcommand{\filesSeedThreeFailed}{5{,}681}
\newcommand{\filesSeedThreeFailedShare}{1.58\%}
\newcommand{\scaleSeedThreeClassRepos}{500}
\newcommand{\scaleSeedThreeXMad}{495}
\newcommand{\scaleSeedThreeDMad}{122}
\newcommand{\scaleSeedThreeDMeanAd}{364}
\newcommand{\scaleSeedThreeDConstant}{14}
\newcommand{\zeroSeedThreeDClass}{65.6\%}
\newcommand{\scaleSeedThreeMethodRepos}{500}
\newcommand{\scaleSeedThreeXmMeanAd}{448}
\newcommand{\scaleSeedThreeDmMeanAd}{492}
\newcommand{\zeroSeedThreeXm}{64.9\%}
\newcommand{\zeroSeedThreeDmMethod}{83.5\%}
\newcommand{\joinFontanaLizardUnits}{376{,}803}
\newcommand{\joinFontanaShare}{99.3\%}
\newcommand{\joinFontanaMeasurerExtra}{4{,}594}
\newcommand{\joinFontanaMeasurerExtraShare}{1.2\%}
\newcommand{\joinFontanaClassUnits}{33{,}118}
\newcommand{\joinFontanaClassCoupledShare}{99.9\%}
\newcommand{\callsFontanaInternal}{49.9\%}
\newcommand{\callsFontanaJdk}{35.7\%}
\newcommand{\callsFontanaUnresolved}{11.1\%}
\newcommand{\callsFontanaShouldHave}{3.2\%}
\newcommand{\callsFontanaBudget}{0.0\%}
\newcommand{\roleFontanaClassCatchmentUnits}{37.9\%}
\newcommand{\roleFontanaClassCatchmentCc}{32.7\%}
\newcommand{\roleFontanaClassPipeUnits}{4.7\%}
\newcommand{\roleFontanaClassPipeCc}{1.1\%}
\newcommand{\roleFontanaClassJunctionUnits}{57.4\%}
\newcommand{\roleFontanaClassJunctionCc}{66.2\%}
\newcommand{\throwNewFontanaClassShare}{1.8\%}
\newcommand{\diFontanaRepos}{3}
\newcommand{\prunedFontanaFiles}{14}
\newcommand{\filesFontana}{37{,}781}
\newcommand{\filesFontanaFailed}{46}
\newcommand{\filesFontanaFailedShare}{0.12\%}
\newcommand{\scaleFontanaClassRepos}{73}
\newcommand{\scaleFontanaXMad}{73}
\newcommand{\scaleFontanaDMad}{54}
\newcommand{\scaleFontanaDMeanAd}{19}
\newcommand{\scaleFontanaDConstant}{0}
\newcommand{\zeroFontanaDClass}{40.0\%}
\newcommand{\uncoupledBlobReviewed}{934}
\newcommand{\uncoupledBlobUncoupled}{0}
\newcommand{\uncoupledLongReviewed}{1034}
\newcommand{\uncoupledLongUncoupled}{7}
\newcommand{\uncoupledLongPositiveUncoupled}{0}
\newcommand{\uncoupledFontanaUncoupled}{3}
\newcommand{\uncoupledFontanaPositiveUncoupled}{2}
\newcommand{\labelsBlobReviewedLabelled}{1{,}060}
\newcommand{\labelsBlobReviewedMeasured}{934}
\newcommand{\labelsBlobLabelled}{110}
\newcommand{\labelsBlobMeasured}{98}
\newcommand{\labelsBlobReviewedDeleted}{5}
\newcommand{\labelsLongReviewedLabelled}{1{,}170}
\newcommand{\labelsLongReviewedMeasured}{1{,}034}
\newcommand{\labelsLongLabelled}{134}
\newcommand{\labelsLongMeasured}{125}
\newcommand{\labelsLongReviewedDeleted}{2}
\newcommand{\mlcqOwnerTopShare}{50.8\%}
\newcommand{\seedThreeOwnerTopShare}{2.8\%}
\newcommand{\mlcqApacheRepos}{131}
\newcommand{\agreeCcExact}{99.7\%}
\newcommand{\agreeCcRho}{0.997}
\newcommand{\agreeCcRhoWithin}{0.998}
\newcommand{\agreeCognitiveExact}{94.6\%}
\newcommand{\agreeCognitiveRho}{0.959}
\newcommand{\agreeCognitiveRhoWithin}{0.972}
\newcommand{\agreeNlocExact}{98.3\%}
\newcommand{\agreeNlocRho}{0.998}
\newcommand{\agreeNlocRhoWithin}{0.999}
\newcommand{\agreeTokensExact}{94.6\%}
\newcommand{\agreeTokensRho}{0.998}
\newcommand{\agreeTokensRhoWithin}{0.999}
\newcommand{\agreeParamsExact}{99.9\%}
\newcommand{\agreeParamsRho}{0.999}
\newcommand{\agreeParamsRhoWithin}{1.000}
\newcommand{\agreeNestingExact}{93.3\%}
\newcommand{\agreeNestingRho}{0.911}
\newcommand{\agreeNestingRhoWithin}{0.939}
\newcommand{\agreeCycDensityExact}{98.3\%}
\newcommand{\agreeCycDensityRho}{0.997}
\newcommand{\agreeCycDensityRhoWithin}{0.998}
\newcommand{\agreeExitPointsExact}{99.1\%}
\newcommand{\agreeExitPointsRho}{0.989}
\newcommand{\agreeExitPointsRhoWithin}{0.994}
\newcommand{\agreeHvolExact}{1.5\%}
\newcommand{\agreeHvolRho}{0.990}
\newcommand{\agreeHvolRhoWithin}{0.990}
\newcommand{\agreeMiExact}{1.5\%}
\newcommand{\agreeMiRho}{0.995}
\newcommand{\agreeMiRhoWithin}{0.994}
\newcommand{\agreeUnits}{2{,}300{,}061}
\newcommand{\agreeRepos}{261}
\newcommand{\agreeJoinShare}{99.4\%}
\newcommand{\pmdCogPaperOne}{0.982}
\newcommand{\pmdCogMeasurer}{0.972}
\newcommand{\callsExploreInternalMedian}{34.8\%}
\newcommand{\callsExploreUnresolvedMedian}{26.9\%}
\newcommand{\callsExploreShouldHaveMedian}{6.8\%}
\newcommand{\xClass}{nloc\_max}
\newcommand{\xClassMcc}{0.287}
\newcommand{\xClassSecond}{hvol\_max}
\newcommand{\xClassSecondMcc}{0.279}
\newcommand{\xClassThird}{tokens\_max}
\newcommand{\xClassThirdMcc}{0.268}
\newcommand{\xClassCogMaxMcc}{0.234}
\newcommand{\xClassCogMaxRank}{6}
\newcommand{\xClassRepos}{205}
\newcommand{\xClassCandidates}{18}
\newcommand{\xMethod}{cognitive}
\newcommand{\xMethodMcc}{0.600}
\newcommand{\xMethodSecond}{cc}
\newcommand{\xMethodSecondMcc}{0.531}
\newcommand{\xMethodRepos}{212}
\newcommand{\gateRequired}{72}
\newcommand{\gateHalfWidth}{0.095}
\newcommand{\gateExploreEstimate}{$-$0.063 [$-$0.155, +0.034]}
\newcommand{\gateExploreRepos}{205}
\newcommand{\gateExplorePositiveRepos}{80}
\newcommand{\gateVOneRequired}{73}
\newcommand{\gateVOneS}{0.434}
\newcommand{\primRepos}{208}
\newcommand{\primReposWithPositive}{66}
\newcommand{\primUniv}{934}
\newcommand{\primPos}{98}
\newcommand{\primNeg}{836}
\newcommand{\primBaseRate}{10.5\%}
\newcommand{\primMicro}{$-$0.041 [$-$0.130, +0.054]}
\newcommand{\primMicroValue}{$-$0.041}
\newcommand{\primMicroLo}{$-$0.130}
\newcommand{\primMicroHi}{+0.054}
\newcommand{\primReading}{no difference resolved}
\newcommand{\primVerdict}{no difference resolved (underpowered)}
\newcommand{\primMacro}{$-$0.004 [$-$0.032, +0.026]}
\newcommand{\primHalfWidthConfirm}{0.092}
\newcommand{\primGateReposWithBlob}{66}
\newcommand{\primXMcc}{0.187}
\newcommand{\primXPrecision}{0.26}
\newcommand{\primXRecall}{0.31}
\newcommand{\primXFpr}{0.104}
\newcommand{\primXTp}{30}
\newcommand{\primXFp}{87}
\newcommand{\primXFn}{68}
\newcommand{\primXTn}{749}
\newcommand{\primXFlagged}{117}
\newcommand{\primDMcc}{0.146}
\newcommand{\primDPrecision}{0.19}
\newcommand{\primDRecall}{0.41}
\newcommand{\primDFpr}{0.208}
\newcommand{\primDTp}{40}
\newcommand{\primDFp}{174}
\newcommand{\primDFn}{58}
\newcommand{\primDTn}{662}
\newcommand{\primDFlagged}{214}
\newcommand{\primAucX}{0.768}
\newcommand{\primAucD}{0.631}
\newcommand{\primAucDiff}{$-$0.137 [$-$0.195, $-$0.080]}
\newcommand{\primJpMicro}{$-$0.050 [$-$0.142, +0.049]}
\newcommand{\primJpVerdict}{no difference resolved (underpowered)}
\newcommand{\primJpUniv}{933}
\newcommand{\primJpAucX}{0.767}
\newcommand{\primJpAucD}{0.631}
\newcommand{\primJpAucDiff}{$-$0.136 [$-$0.194, $-$0.079]}
\newcommand{\primCtorMicro}{$-$0.062 [$-$0.130, +0.011]}
\newcommand{\primCtorReading}{no difference resolved}
\newcommand{\primCtorAgrees}{agrees}
\newcommand{\primCtorMcc}{0.125}
\newcommand{\primCtorPrecision}{0.20}
\newcommand{\primCtorRecall}{0.27}
\newcommand{\primMethodsVsCtorMicro}{+0.021 [$-$0.048, +0.097]}
\newcommand{\primEMicro}{+0.025 [$-$0.052, +0.100]}
\newcommand{\primEReading}{no difference resolved}
\newcommand{\primEMacro}{+0.011 [$-$0.014, +0.038]}
\newcommand{\primEYMcc}{0.212}
\newcommand{\primEYPrecision}{0.26}
\newcommand{\primEYRecall}{0.37}
\newcommand{\primEYFp}{102}
\newcommand{\primEYTp}{36}
\newcommand{\primAmpMicro}{+0.009 [$-$0.075, +0.104]}
\newcommand{\primAmpReading}{no difference resolved}
\newcommand{\primAmpMacro}{+0.019 [$-$0.013, +0.056]}
\newcommand{\primAmpYMcc}{0.196}
\newcommand{\primAmpYPrecision}{0.20}
\newcommand{\primAmpYRecall}{0.56}
\newcommand{\primAmpYFp}{224}
\newcommand{\primAmpYTp}{55}
\newcommand{\primDvsEMicro}{$-$0.066 [$-$0.130, $-$0.000]}
\newcommand{\primDvsEReading}{exemption only beats the whole factor}
\newcommand{\primDvsEMacro}{$-$0.015 [$-$0.038, +0.008]}
\newcommand{\primDvsEYMcc}{0.146}
\newcommand{\primDvsEYPrecision}{0.19}
\newcommand{\primDvsEYRecall}{0.41}
\newcommand{\primDvsEYFp}{174}
\newcommand{\primDvsEYTp}{40}
\newcommand{\primDvsAMicro}{$-$0.051 [$-$0.112, +0.006]}
\newcommand{\primDvsAReading}{no difference resolved}
\newcommand{\primDvsAMacro}{$-$0.023 [$-$0.046, $-$0.004]}
\newcommand{\primDvsAYMcc}{0.146}
\newcommand{\primDvsAYPrecision}{0.19}
\newcommand{\primDvsAYRecall}{0.41}
\newcommand{\primDvsAYFp}{174}
\newcommand{\primDvsAYTp}{40}
\newcommand{\primNvsXMicro}{+0.049 [$-$0.023, +0.121]}
\newcommand{\primNvsXReading}{no difference resolved}
\newcommand{\primNvsXMacro}{+0.029 [0.000, +0.058]}
\newcommand{\primDvsNMicro}{$-$0.091 [$-$0.175, $-$0.006]}
\newcommand{\primDvsNReading}{the size factor beats the coupling factor}
\newcommand{\primDvsNMacro}{$-$0.033 [$-$0.072, +0.004]}
\newcommand{\primNMcc}{0.236}
\newcommand{\primNPrecision}{0.22}
\newcommand{\primNRecall}{0.60}
\newcommand{\primNTp}{59}
\newcommand{\primNFp}{211}
\newcommand{\primMechanism}{not read: the primary does not read refinement}
\newcommand{\primJpDvsAReading}{amplification only beats the whole factor}
\newcommand{\primJpDvsAMicro}{$-$0.062 [$-$0.129, $-$0.004]}
\newcommand{\primJpDvsEMicro}{$-$0.078 [$-$0.148, $-$0.008]}
\newcommand{\primJpDvsNMicro}{$-$0.097 [$-$0.184, $-$0.009]}
\newcommand{\primJpEMicro}{+0.028 [$-$0.048, +0.102]}
\newcommand{\primJpAmpMicro}{+0.012 [$-$0.074, +0.105]}
\newcommand{\leafCatchFlagged}{38}
\newcommand{\leafCatchPositive}{8}
\newcommand{\leafCatchRate}{21.1\%}
\newcommand{\leafRestFlagged}{79}
\newcommand{\leafRestPositive}{22}
\newcommand{\leafRestRate}{27.8\%}
\newcommand{\leafDiff}{$-$0.068 [$-$0.238, +0.102]}
\newcommand{\leafReading}{unresolved}
\newcommand{\leafMethodCatchFlagged}{28}
\newcommand{\leafMethodCatchPositive}{13}
\newcommand{\leafMethodCatchRate}{46.4\%}
\newcommand{\leafMethodRestRate}{76.8\%}
\newcommand{\leafMethodDiff}{$-$0.304 [$-$0.541, $-$0.082]}
\newcommand{\leafClassUnresolved}{12}
\newcommand{\leafClassVariants}{12}
\newcommand{\leafMethodReading}{supported}
\newcommand{\leafMethodSupported}{3}
\newcommand{\leafMethodVariants}{3}
\newcommand{\leafFontanaReading}{supported}
\newcommand{\leafFontanaSupported}{8}
\newcommand{\leafFontanaFlowSupported}{4}
\newcommand{\leafFontanaVariants}{12}
\newcommand{\leafMethodDataDiff}{$-$0.208 [$-$0.413, $-$0.017]}
\newcommand{\leafMethodMessageDiff}{$-$0.327 [$-$0.557, $-$0.102]}
\newcommand{\leafFontanaCatchRate}{75.0\%}
\newcommand{\leafFontanaRestRate}{94.9\%}
\newcommand{\leafFontanaDiff}{$-$0.199 [$-$0.362, $-$0.001]}
\newcommand{\leafFontanaCatchFlagged}{20}
\newcommand{\leafFontanaCatchPositive}{15}
\newcommand{\consVerdict}{outgoing data only}
\newcommand{\consCtorVerdict}{outgoing data only}
\newcommand{\consCtorAgrees}{agrees}
\newcommand{\consJpVerdict}{outgoing data only}
\newcommand{\consJpAgrees}{agrees}
\newcommand{\pOneHolds}{does not hold}
\newcommand{\pOneDPeak}{+0.041 [$-$0.031, +0.114]}
\newcommand{\pOneDDensity}{+0.094 [+0.013, +0.174]}
\newcommand{\pOneOPeak}{+0.125 [+0.054, +0.195]}
\newcommand{\pOneODensity}{+0.151 [+0.076, +0.227]}
\newcommand{\pOneOHolds}{holds}
\newcommand{\pTwoHolds}{holds}
\newcommand{\pTwoBlob}{+0.073 [+0.024, +0.122]}
\newcommand{\pTwoLong}{+0.188 [+0.109, +0.266]}
\newcommand{\pTwoFontana}{+0.242 [+0.195, +0.295]}
\newcommand{\pTwoMinToolLo}{+0.013}
\newcommand{\pTwoJpBlob}{+0.065 [+0.019, +0.111]}
\newcommand{\pTwoJpHolds}{holds}
\newcommand{\pThreeReads}{outgoing-only}
\newcommand{\pThreeSize}{$-$0.080 [$-$0.144, $-$0.014]}
\newcommand{\pThreePeak}{$-$0.084 [$-$0.139, $-$0.029]}
\newcommand{\pThreeDensity}{$-$0.058 [$-$0.110, $-$0.009]}
\newcommand{\pFourPositive}{5}
\newcommand{\pFourCells}{12}
\newcommand{\pFourMajority}{not a majority}
\newcommand{\pFourLoAbove}{5}
\newcommand{\pFourHiBelow}{4}
\newcommand{\pOneDataHolds}{does not hold}
\newcommand{\pOneDataOHolds}{holds}
\newcommand{\pThreeDataReads}{outgoing-only}
\newcommand{\pOneDataDPeak}{+0.061 [$-$0.014, +0.137]}
\newcommand{\pOneDataDDensity}{+0.109 [+0.024, +0.192]}
\newcommand{\pOneDataOPeak}{+0.103 [+0.033, +0.176]}
\newcommand{\pOneDataODensity}{+0.131 [+0.054, +0.204]}
\newcommand{\pOneMessageHolds}{does not hold}
\newcommand{\pOneMessageOHolds}{holds}
\newcommand{\pThreeMessageReads}{outgoing-only}
\newcommand{\pOneMessageDPeak}{+0.066 [$-$0.001, +0.145]}
\newcommand{\pOneMessageDDensity}{+0.107 [+0.026, +0.189]}
\newcommand{\pOneMessageOPeak}{+0.130 [+0.059, +0.199]}
\newcommand{\pOneMessageODensity}{+0.147 [+0.072, +0.219]}
\newcommand{\cellBlobSizeD}{$-$0.032 [$-$0.114, +0.048]}
\newcommand{\cellBlobSizeO}{+0.048 [$-$0.014, +0.108]}
\newcommand{\cellBlobSizeDO}{$-$0.080 [$-$0.144, $-$0.014]}
\newcommand{\cellBlobSizeN}{208}
\newcommand{\cellBlobPeakD}{+0.041 [$-$0.031, +0.114]}
\newcommand{\cellBlobPeakO}{+0.125 [+0.054, +0.195]}
\newcommand{\cellBlobPeakDO}{$-$0.084 [$-$0.139, $-$0.029]}
\newcommand{\cellBlobPeakN}{208}
\newcommand{\cellBlobDensityD}{+0.094 [+0.013, +0.174]}
\newcommand{\cellBlobDensityO}{+0.151 [+0.076, +0.227]}
\newcommand{\cellBlobDensityDO}{$-$0.058 [$-$0.110, $-$0.009]}
\newcommand{\cellBlobDensityN}{208}
\newcommand{\cellLongSizeD}{$-$0.232 [$-$0.317, $-$0.144]}
\newcommand{\cellLongSizeO}{$-$0.159 [$-$0.200, $-$0.121]}
\newcommand{\cellLongSizeDO}{$-$0.073 [$-$0.156, +0.012]}
\newcommand{\cellLongSizeN}{202}
\newcommand{\cellLongDensityD}{+0.280 [+0.197, +0.362]}
\newcommand{\cellLongDensityO}{+0.222 [+0.162, +0.280]}
\newcommand{\cellLongDensityDO}{+0.058 [+0.006, +0.110]}
\newcommand{\cellLongDensityN}{202}
\newcommand{\cellLongStructureD}{$-$0.044 [$-$0.115, +0.028]}
\newcommand{\cellLongStructureO}{$-$0.028 [$-$0.086, +0.034]}
\newcommand{\cellLongStructureDO}{$-$0.017 [$-$0.061, +0.024]}
\newcommand{\cellLongStructureN}{202}
\newcommand{\cellFontanaSizeD}{$-$0.073 [$-$0.193, +0.037]}
\newcommand{\cellFontanaSizeO}{+0.063 [$-$0.004, +0.121]}
\newcommand{\cellFontanaSizeDO}{$-$0.136 [$-$0.243, $-$0.038]}
\newcommand{\cellFontanaSizeN}{72}
\newcommand{\cellFontanaPeakD}{+0.169 [+0.068, +0.269]}
\newcommand{\cellFontanaPeakO}{+0.315 [+0.249, +0.381]}
\newcommand{\cellFontanaPeakDO}{$-$0.146 [$-$0.230, $-$0.067]}
\newcommand{\cellFontanaPeakN}{72}
\newcommand{\cellFontanaDensityD}{+0.413 [+0.326, +0.503]}
\newcommand{\cellFontanaDensityO}{+0.573 [+0.499, +0.658]}
\newcommand{\cellFontanaDensityDO}{$-$0.160 [$-$0.247, $-$0.079]}
\newcommand{\cellFontanaDensityN}{72}
\newcommand{\cellsCountedFlow}{21}
\newcommand{\cellsCountedDLoAbove}{9}
\newcommand{\cellsCountedDHiBelow}{5}
\newcommand{\cellsCountedOLoAbove}{12}
\newcommand{\cellsCountedOHiBelow}{6}
\newcommand{\cellsCountedDOHiBelow}{12}
\newcommand{\cellsCountedDOLoAbove}{5}
\newcommand{\cellsCountedClassFlow}{12}
\newcommand{\cellsCountedClassDOHiBelow}{12}
\newcommand{\cellsDensityFlow}{7}
\newcommand{\cellsDensityDLoAbove}{7}
\newcommand{\cellsSizeFlow}{7}
\newcommand{\cellsSizeDHiBelow}{5}
\newcommand{\betweenCells}{34}
\newcommand{\betweenFlowDataLoAbove}{10}
\newcommand{\betweenFlowDataHiBelow}{1}
\newcommand{\betweenMessageFlowLoAbove}{2}
\newcommand{\betweenMessageFlowHiBelow}{0}
\newcommand{\betweenBlobPeakMessageFlow}{+0.025 [+0.004, +0.050]}
\newcommand{\betweenLongStructureMessageFlow}{+0.019 [$-$0.004, +0.042]}
\newcommand{\betweenFontanaDensityMessageFlow}{+0.045 [+0.003, +0.100]}
\newcommand{\betweenBlobPeakFlowData}{$-$0.020 [$-$0.054, +0.014]}
\newcommand{\secDDataMicro}{$-$0.008 [$-$0.101, +0.096]}
\newcommand{\secDDataMacro}{+0.013 [$-$0.018, +0.048]}
\newcommand{\secDDataReading}{no difference resolved}
\newcommand{\secDDataMcc}{0.180}
\newcommand{\secDDataPrecision}{0.20}
\newcommand{\secDDataRecall}{0.47}
\newcommand{\secDDataFp}{182}
\newcommand{\secDDataTp}{46}
\newcommand{\secDMessageMicro}{$-$0.005 [$-$0.089, +0.084]}
\newcommand{\secDMessageMacro}{+0.012 [$-$0.017, +0.044]}
\newcommand{\secDMessageReading}{no difference resolved}
\newcommand{\secDMessageMcc}{0.182}
\newcommand{\secDMessagePrecision}{0.20}
\newcommand{\secDMessageRecall}{0.48}
\newcommand{\secDMessageFp}{186}
\newcommand{\secDMessageTp}{47}
\newcommand{\secOFlowMicro}{+0.085 [+0.001, +0.182]}
\newcommand{\secOFlowMacro}{+0.063 [+0.028, +0.103]}
\newcommand{\secOFlowReading}{the divergence factor refines complexity}
\newcommand{\secOFlowMcc}{0.272}
\newcommand{\secOFlowPrecision}{0.23}
\newcommand{\secOFlowRecall}{0.69}
\newcommand{\secOFlowFp}{233}
\newcommand{\secOFlowTp}{68}
\newcommand{\secODataMicro}{+0.035 [$-$0.053, +0.134]}
\newcommand{\secODataMacro}{+0.040 [+0.006, +0.076]}
\newcommand{\secODataReading}{no difference resolved}
\newcommand{\secODataMcc}{0.222}
\newcommand{\secODataPrecision}{0.21}
\newcommand{\secODataRecall}{0.61}
\newcommand{\secODataFp}{231}
\newcommand{\secODataTp}{60}
\newcommand{\secOMessageMicro}{+0.087 [+0.003, +0.183]}
\newcommand{\secOMessageMacro}{+0.058 [+0.025, +0.097]}
\newcommand{\secOMessageReading}{the divergence factor refines complexity}
\newcommand{\secOMessageMcc}{0.274}
\newcommand{\secOMessagePrecision}{0.23}
\newcommand{\secOMessageRecall}{0.68}
\newcommand{\secOMessageFp}{225}
\newcommand{\secOMessageTp}{67}
\newcommand{\secDAloneMicro}{$-$0.014 [$-$0.109, +0.091]}
\newcommand{\secDAloneMacro}{+0.004 [$-$0.027, +0.038]}
\newcommand{\secDAloneReading}{no difference resolved}
\newcommand{\secDAloneMcc}{0.173}
\newcommand{\secDAlonePrecision}{0.21}
\newcommand{\secDAloneRecall}{0.40}
\newcommand{\secDAloneFp}{145}
\newcommand{\secDAloneTp}{39}
\newcommand{\secOAloneMicro}{0.000 [$-$0.100, +0.106]}
\newcommand{\secOAloneMacro}{+0.019 [$-$0.017, +0.058]}
\newcommand{\secOAloneReading}{no difference resolved}
\newcommand{\secOAloneMcc}{0.187}
\newcommand{\secOAlonePrecision}{0.21}
\newcommand{\secOAloneRecall}{0.45}
\newcommand{\secOAloneFp}{163}
\newcommand{\secOAloneTp}{44}
\newcommand{\secAndMicro}{$-$0.023 [$-$0.092, +0.045]}
\newcommand{\secAndMacro}{$-$0.015 [$-$0.037, +0.004]}
\newcommand{\secAndReading}{no difference resolved}
\newcommand{\secAndMcc}{0.164}
\newcommand{\secAndPrecision}{0.30}
\newcommand{\secAndRecall}{0.17}
\newcommand{\secAndFp}{39}
\newcommand{\secAndTp}{17}
\newcommand{\secDPubMicro}{$-$0.027 [$-$0.116, +0.071]}
\newcommand{\secDPubMacro}{+0.003 [$-$0.025, +0.034]}
\newcommand{\secDPubReading}{no difference resolved}
\newcommand{\secDPubMcc}{0.161}
\newcommand{\secDPubPrecision}{0.18}
\newcommand{\secDPubRecall}{0.49}
\newcommand{\secDPubFp}{213}
\newcommand{\secDPubTp}{48}
\newcommand{\secDFormBMicro}{+0.044 [$-$0.038, +0.139]}
\newcommand{\secDFormBMacro}{+0.035 [+0.003, +0.073]}
\newcommand{\secDFormBReading}{no difference resolved}
\newcommand{\secDFormBMcc}{0.231}
\newcommand{\secDFormBPrecision}{0.21}
\newcommand{\secDFormBRecall}{0.60}
\newcommand{\secDFormBFp}{216}
\newcommand{\secDFormBTp}{59}
\newcommand{\secDFormBCtorMicro}{+0.064 [$-$0.023, +0.162]}
\newcommand{\secDFormBCtorMacro}{+0.047 [+0.014, +0.087]}
\newcommand{\secDFormBCtorReading}{no difference resolved}
\newcommand{\secDFormBCtorMcc}{0.251}
\newcommand{\secDFormBCtorPrecision}{0.22}
\newcommand{\secDFormBCtorRecall}{0.63}
\newcommand{\secDFormBCtorFp}{216}
\newcommand{\secDFormBCtorTp}{62}
\newcommand{\secOFormBMicro}{+0.026 [$-$0.046, +0.107]}
\newcommand{\secOFormBMacro}{+0.027 [$-$0.004, +0.061]}
\newcommand{\secOFormBReading}{no difference resolved}
\newcommand{\secOFormBMcc}{0.213}
\newcommand{\secOFormBPrecision}{0.21}
\newcommand{\secOFormBRecall}{0.54}
\newcommand{\secOFormBFp}{195}
\newcommand{\secOFormBTp}{53}
\newcommand{\secOtherArms}{41}
\newcommand{\secOtherLoAbove}{2}
\newcommand{\secOtherHiBelow}{0}
\newcommand{\secOtherMacroLoAbove}{12}
\newcommand{\secFormBArms}{15}
\newcommand{\secFormBMacroLoAbove}{9}
\newcommand{\secFormBMinMicro}{$-$0.003}
\newcommand{\secFormBMaxMicro}{+0.064}
\newcommand{\rhoXD}{0.21}
\newcommand{\rhoXOut}{0.55}
\newcommand{\rhoXDData}{0.10}
\newcommand{\rhoNomD}{$-$0.23}
\newcommand{\rhoNomOut}{0.40}
\newcommand{\rhoNlocSumOut}{0.57}
\newcommand{\rhoXDSeed}{0.16}
\newcommand{\rhoXOutSeed}{0.52}
\newcommand{\rhoXmD}{0.25}
\newcommand{\rhoXmOut}{0.39}
\newcommand{\trackClosely}{do not track closely}
\newcommand{\tertileOne}{$-$0.012 [$-$0.021, 0.000]}
\newcommand{\tertileTwo}{$-$0.026 [$-$0.115, +0.085]}
\newcommand{\tertileThree}{$-$0.022 [$-$0.151, +0.118]}
\newcommand{\methRepos}{202}
\newcommand{\methUniv}{1033}
\newcommand{\methPos}{125}
\newcommand{\methBaseRate}{12.1\%}
\newcommand{\methXMcc}{0.603}
\newcommand{\methXPrecision}{0.69}
\newcommand{\methXRecall}{0.61}
\newcommand{\methXTp}{76}
\newcommand{\methXFp}{34}
\newcommand{\methXFn}{49}
\newcommand{\methDMcc}{0.511}
\newcommand{\methDPrecision}{0.86}
\newcommand{\methDRecall}{0.34}
\newcommand{\methDTp}{43}
\newcommand{\methDFp}{7}
\newcommand{\methDFn}{82}
\newcommand{\methMicro}{$-$0.092 [$-$0.188, +0.010]}
\newcommand{\methReading}{no difference resolved}
\newcommand{\methMacro}{$-$0.057 [$-$0.099, $-$0.015]}
\newcommand{\methNoHelpMicro}{$-$0.144 [$-$0.244, $-$0.043]}
\newcommand{\methNoHelpReading}{it makes it worse}
\newcommand{\methNoHelpMcc}{0.460}
\newcommand{\methNoHelpVsAllMicro}{$-$0.051 [$-$0.125, +0.016]}
\newcommand{\methEMicro}{$-$0.002 [$-$0.077, +0.073]}
\newcommand{\methEReading}{no difference resolved}
\newcommand{\methEMcc}{0.601}
\newcommand{\methERecall}{0.58}
\newcommand{\methAmpMicro}{$-$0.093 [$-$0.174, $-$0.007]}
\newcommand{\methAmpReading}{the amplification makes it worse}
\newcommand{\methAmpMcc}{0.510}
\newcommand{\methAmpRecall}{0.40}
\newcommand{\methDvsEMicro}{$-$0.090 [$-$0.173, $-$0.006]}
\newcommand{\methDvsEReading}{exemption only beats the whole factor}
\newcommand{\methDvsAMicro}{+0.001 [$-$0.043, +0.040]}
\newcommand{\methDvsAReading}{no difference resolved}
\newcommand{\methNeVsXMicro}{$-$0.060 [$-$0.153, +0.045]}
\newcommand{\methNeVsXReading}{no difference resolved}
\newcommand{\methNeVsAllMicro}{+0.032 [$-$0.058, +0.126]}
\newcommand{\methNeVsAllMacro}{+0.037 [+0.004, +0.071]}
\newcommand{\methNeMcc}{0.543}
\newcommand{\methNePrecision}{0.54}
\newcommand{\methNeRecall}{0.68}
\newcommand{\methMechanism}{not read: the primary does not read refinement}
\newcommand{\methJpMicro}{$-$0.079 [$-$0.171, +0.026]}
\newcommand{\methJpReading}{no difference resolved}
\newcommand{\toolPmdMicro}{$-$0.143 [$-$0.176, $-$0.106]}
\newcommand{\toolPmdReading}{it makes it worse}
\newcommand{\toolPmdN}{490}
\newcommand{\toolPmdXMcc}{0.351}
\newcommand{\toolPmdDMcc}{0.208}
\newcommand{\toolPmdXRecall}{0.43}
\newcommand{\toolPmdDRecall}{0.52}
\newcommand{\toolPmdXPrecision}{0.37}
\newcommand{\toolPmdDPrecision}{0.17}
\newcommand{\toolPmdCtorMicro}{$-$0.117 [$-$0.153, $-$0.078]}
\newcommand{\toolDecorMicro}{$-$0.093 [$-$0.121, $-$0.059]}
\newcommand{\toolDecorReading}{it makes it worse}
\newcommand{\toolDecorN}{490}
\newcommand{\toolDecorXMcc}{0.202}
\newcommand{\toolDecorDMcc}{0.109}
\newcommand{\toolDecorXRecall}{0.25}
\newcommand{\toolDecorDRecall}{0.35}
\newcommand{\toolDecorXPrecision}{0.29}
\newcommand{\toolDecorDPrecision}{0.15}
\newcommand{\toolDecorCtorMicro}{$-$0.051 [$-$0.097, $-$0.002]}
\newcommand{\toolWmcMicro}{$-$0.144 [$-$0.185, $-$0.099]}
\newcommand{\toolWmcReading}{it makes it worse}
\newcommand{\toolWmcN}{491}
\newcommand{\toolWmcXMcc}{0.390}
\newcommand{\toolWmcDMcc}{0.247}
\newcommand{\toolWmcXRecall}{0.38}
\newcommand{\toolWmcDRecall}{0.50}
\newcommand{\toolWmcXPrecision}{0.52}
\newcommand{\toolWmcDPrecision}{0.25}
\newcommand{\toolWmcCtorMicro}{$-$0.122 [$-$0.161, $-$0.074]}
\newcommand{\toolNpathMicro}{$-$0.021 [$-$0.054, +0.013]}
\newcommand{\toolNpathReading}{no difference resolved}
\newcommand{\toolNpathN}{490}
\newcommand{\toolNpathXMcc}{0.350}
\newcommand{\toolNpathDMcc}{0.330}
\newcommand{\toolNpathXRecall}{0.97}
\newcommand{\toolNpathDRecall}{0.62}
\newcommand{\toolNpathXPrecision}{0.13}
\newcommand{\toolNpathDPrecision}{0.19}
\newcommand{\toolBrainMicro}{+0.036 [+0.025, +0.047]}
\newcommand{\toolBrainReading}{the divergence factor refines complexity}
\newcommand{\toolBrainN}{491}
\newcommand{\toolBrainXMcc}{0.220}
\newcommand{\toolBrainDMcc}{0.256}
\newcommand{\toolBrainXRecall}{1.00}
\newcommand{\toolBrainDRecall}{0.78}
\newcommand{\toolBrainXPrecision}{0.05}
\newcommand{\toolBrainDPrecision}{0.09}
\newcommand{\toolBrainJpMicro}{+0.042 [+0.032, +0.053]}
\newcommand{\fontanaPopSystems}{72}
\newcommand{\fontanaPopUniv}{380}
\newcommand{\fontanaPopPos}{139}
\newcommand{\fontanaMicro}{$-$0.093 [$-$0.220, +0.035]}
\newcommand{\fontanaReading}{no difference resolved}
\newcommand{\fontanaMacro}{$-$0.065 [$-$0.155, +0.021]}
\newcommand{\fontanaXMcc}{0.765}
\newcommand{\fontanaDMcc}{0.672}
\newcommand{\fontanaXPrecision}{0.91}
\newcommand{\fontanaDPrecision}{0.79}
\newcommand{\fontanaXRecall}{0.78}
\newcommand{\fontanaDRecall}{0.80}
\newcommand{\fontanaXFp}{11}
\newcommand{\fontanaDFp}{30}
\newcommand{\fontanaCtorMicro}{$-$0.065 [$-$0.169, +0.039]}
\newcommand{\fontanaCtorReading}{no difference resolved}
\newcommand{\fontanaCtorMacro}{$-$0.084 [$-$0.167, $-$0.004]}
\newcommand{\fontanaVOneSystems}{19}
\newcommand{\fontanaVOneUniv}{100}
\newcommand{\fontanaVOnePos}{46}
\newcommand{\fontanaVOneMicro}{$-$0.203 [$-$0.337, $-$0.074]}
\newcommand{\fontanaVOneReading}{it makes it worse}
\newcommand{\fontanaVOneMeasured}{46}
\newcommand{\primVOneMicro}{$-$0.075 [$-$0.191, +0.045]}
\newcommand{\consVOneVerdict}{outgoing data only}
\newcommand{\fontanaPmdMedianMeasurer}{0.978}
\newcommand{\cellFontanaDensityOFlag}{+0.573 [+0.499, +0.658]}
\newcommand{\friedBlobSizeRankX}{2.08}
\newcommand{\friedBlobSizeRankD}{2.16}
\newcommand{\friedBlobSizeRankO}{1.76}
\newcommand{\friedBlobSizeP}{$p = 0.100$}
\newcommand{\friedBlobSizeN}{53}
\newcommand{\friedBlobPeakRankX}{2.32}
\newcommand{\friedBlobPeakRankD}{2.11}
\newcommand{\friedBlobPeakRankO}{1.58}
\newcommand{\friedBlobPeakP}{$p < 0.001$}
\newcommand{\friedBlobPeakN}{52}
\newcommand{\friedBlobDensityRankX}{2.37}
\newcommand{\friedBlobDensityRankD}{2.00}
\newcommand{\friedBlobDensityRankO}{1.63}
\newcommand{\friedBlobDensityP}{$p < 0.001$}
\newcommand{\friedBlobDensityN}{53}
\newcommand{\friedLongSizeRankX}{1.40}
\newcommand{\friedLongSizeRankD}{2.40}
\newcommand{\friedLongSizeRankO}{2.20}
\newcommand{\friedLongSizeP}{$p < 0.001$}
\newcommand{\friedLongSizeN}{63}
\newcommand{\friedLongDensityRankX}{2.55}
\newcommand{\friedLongDensityRankD}{1.71}
\newcommand{\friedLongDensityRankO}{1.74}
\newcommand{\friedLongDensityP}{$p < 0.001$}
\newcommand{\friedLongDensityN}{66}
\newcommand{\friedLongStructureRankX}{1.95}
\newcommand{\friedLongStructureRankD}{2.14}
\newcommand{\friedLongStructureRankO}{1.92}
\newcommand{\friedLongStructureP}{$p = 0.387$}
\newcommand{\friedLongStructureN}{65}
\newcommand{\friedFontanaSizeRankX}{1.99}
\newcommand{\friedFontanaSizeRankD}{2.09}
\newcommand{\friedFontanaSizeRankO}{1.92}
\newcommand{\friedFontanaSizeP}{$p = 0.723$}
\newcommand{\friedFontanaSizeN}{44}
\newcommand{\friedFontanaPeakRankX}{2.65}
\newcommand{\friedFontanaPeakRankD}{1.85}
\newcommand{\friedFontanaPeakRankO}{1.50}
\newcommand{\friedFontanaPeakP}{$p < 0.001$}
\newcommand{\friedFontanaPeakN}{48}
\newcommand{\friedFontanaDensityRankX}{2.84}
\newcommand{\friedFontanaDensityRankD}{1.85}
\newcommand{\friedFontanaDensityRankO}{1.30}
\newcommand{\friedFontanaDensityP}{$p < 0.001$}
\newcommand{\friedFontanaDensityN}{48}
\newcommand{\matrixBlobRows}{16{,}450}
\newcommand{\matrixLongRows}{2{,}926}
\newcommand{\sensDiMicro}{$-$0.108 [$-$0.244, +0.036]}
\newcommand{\sensDiReading}{no difference resolved}
\newcommand{\sensDiRepos}{116}
\newcommand{\sensDiXMcc}{0.225}
\newcommand{\sensDiDMcc}{0.117}
\newcommand{\sensDiUniv}{441}
\newcommand{\sensNeVsXMicro}{+0.017 [$-$0.072, +0.117]}
\newcommand{\sensNeVsXReading}{no difference resolved}
\newcommand{\sensNeVsAllMicro}{+0.058 [+0.001, +0.117]}
\newcommand{\sensNeVsAllReading}{fitting on the non-exempt units is better}
\newcommand{\sensNeMcc}{0.204}
\newcommand{\sensNePrecision}{0.25}
\newcommand{\sensNeRecall}{0.37}
\newcommand{\sensNeFp}{107}
\newcommand{\sensJpNeVsAllMicro}{+0.068 [+0.006, +0.134]}
\newcommand{\sensJpNeVsAllReading}{fitting on the non-exempt units is better}
\newcommand{\sensUncoupledMicro}{$-$0.041 [$-$0.130, +0.054]}
\newcommand{\sensBudgetRepos}{0}
\newcommand{\interBlobBothTp}{17}
\newcommand{\interBlobBothFp}{39}
\newcommand{\interBlobBothPrecision}{30.4\%}
\newcommand{\interBlobXOnlyPrecision}{21.3\%}
\newcommand{\interBlobDOnlyPrecision}{21.7\%}
\newcommand{\interBlobEitherMcc}{0.223}
\newcommand{\interBlobBothVsX}{$-$0.023 [$-$0.092, +0.045]}
\newcommand{\interBlobEitherVsX}{+0.036 [$-$0.023, +0.106]}
\newcommand{\interBlobDVsX}{+0.007 [$-$0.080, +0.104]}
\newcommand{\interBlobDMcc}{0.194}
\newcommand{\interLongBothTp}{47}
\newcommand{\interLongBothFp}{7}
\newcommand{\interLongBothPrecision}{87.0\%}
\newcommand{\interLongDOnlyPrecision}{27.3\%}
\newcommand{\interLongBothVsX}{$-$0.064 [$-$0.147, +0.015]}
\newcommand{\interLongEitherVsX}{$-$0.028 [$-$0.091, +0.043]}
\newcommand{\interLongDVsX}{$-$0.139 [$-$0.238, $-$0.032]}
\newcommand{\interLongDMcc}{0.464}
\newcommand{\interFontanaEitherVsX}{+0.052 [+0.011, +0.094]}
\newcommand{\interFontanaEitherMcc}{0.818}
\newcommand{\interFontanaBothVsX}{$-$0.066 [$-$0.126, $-$0.005]}
\newcommand{\interFontanaBothPrecision}{98.8\%}
\newcommand{\primJpReposWithPositive}{66}
\newcommand{\analysisVTwoShaShort}{9519e2466f44}
\newcommand{\priorMadeUtc}{2026-10-26 09:49:15 UTC}
\newcommand{\priorPrimRepos}{67}
\newcommand{\priorPrimUniv}{548}
\newcommand{\priorPrimPos}{98}
\newcommand{\priorPrimMicro}{$-$0.075 [$-$0.191, +0.045]}
\newcommand{\priorPrimVerdict}{no difference resolved (underpowered)}
\newcommand{\priorPrimXMcc}{0.212}
\newcommand{\priorPrimDMcc}{0.136}
\newcommand{\priorPrimXPrecision}{0.38}
\newcommand{\priorPrimDPrecision}{0.26}
\newcommand{\priorPrimAucX}{0.767}
\newcommand{\priorPrimAucD}{0.620}
\newcommand{\priorPrimAucDiff}{$-$0.147 [$-$0.201, $-$0.093]}
\newcommand{\priorDvsNMicro}{$-$0.127 [$-$0.230, $-$0.021]}
\newcommand{\priorDvsEMicro}{$-$0.087 [$-$0.166, $-$0.012]}
\newcommand{\priorGateReposWithBlob}{66}
\newcommand{\priorGateRequired}{108}
\newcommand{\priorLeafDiff}{$-$0.016 [$-$0.273, +0.271]}
\newcommand{\priorLeafCatchRate}{36.4\%}
\newcommand{\priorLeafRestRate}{37.9\%}
\newcommand{\priorConsVerdict}{outgoing data only}
\newcommand{\priorPOneDPeak}{+0.018 [$-$0.068, +0.112]}
\newcommand{\priorPOneDDensity}{+0.062 [$-$0.041, +0.161]}
\newcommand{\priorPOneOPeak}{+0.147 [+0.061, +0.234]}
\newcommand{\priorPOneODensity}{+0.162 [+0.068, +0.255]}
\newcommand{\priorPThreePeak}{$-$0.128 [$-$0.199, $-$0.062]}
\newcommand{\priorMethRepos}{78}
\newcommand{\priorMethUniv}{711}
\newcommand{\priorMethPos}{125}
\newcommand{\priorMethMicro}{$-$0.136 [$-$0.231, $-$0.041]}
\newcommand{\priorMethReading}{it makes it worse}
\newcommand{\priorFontanaSystems}{51}
\newcommand{\priorFontanaUniv}{305}
\newcommand{\priorFontanaMicro}{$-$0.118 [$-$0.260, +0.020]}
\newcommand{\priorFontanaReading}{no difference resolved}
\newcommand{\priorJpPrimMicro}{$-$0.084 [$-$0.205, +0.032]}
\newcommand{\restoredBlobRepos}{141}
\newcommand{\restoredBlobUnits}{386}
\newcommand{\restoredLongRepos}{124}
\newcommand{\restoredLongUnits}{322}
\newcommand{\restoredFontanaSystems}{21}
\newcommand{\restoredFontanaUnits}{75}
\newcommand{\priorXClass}{nloc\_max}
\newcommand{\priorXClassMcc}{0.277}
\newcommand{\priorXClassRepos}{81}
\newcommand{\priorXMethod}{cognitive}
\newcommand{\priorXMethodMcc}{0.638}
\newcommand{\priorXMethodRepos}{69}
\newcommand{\priorGateRequiredExplore}{108}
\newcommand{\priorGateHalfWidth}{0.116}
\newcommand{\priorGateExploreEstimate}{$-$0.063 [$-$0.178, +0.054]}
\newcommand{\priorGateExploreRepos}{81}
\newcommand{\priorDecisionsMadeUtc}{2026-10-26 07:54:14 UTC}
\newcommand{\restoredExploreBlobRepos}{124}
\newcommand{\restoredExploreLongRepos}{143}
\newcommand{\xClassChanged}{unchanged}
\newcommand{\xMethodChanged}{unchanged}
\newcommand{\expMadeUtc}{20261026 232945 UTC}
\newcommand{\expShaShort}{4bd479a1814c}
\newcommand{\expClassForms}{24}
\newcommand{\expMethodForms}{17}
\newcommand{\expBlobConfirmRepos}{208}
\newcommand{\expBlobExploreRepos}{205}
\newcommand{\expLongConfirmRepos}{202}
\newcommand{\expLongExploreRepos}{212}
\newcommand{\expClassTopConfirm}{lines summed $\times$ (1 + leaving calls)}
\newcommand{\expClassTopConfirmMcc}{0.286}
\newcommand{\expClassTopExplore}{lines summed $\times$ (1 + all internal calls)}
\newcommand{\expClassTopExploreMcc}{0.325}
\newcommand{\expSizeCallsConfirmMcc}{0.286}
\newcommand{\expSizeCallsExploreMcc}{0.314}
\newcommand{\expSizeCallsFontanaMcc}{0.822}
\newcommand{\expSizeCallsConfirmVsComplexity}{+0.032 [$-$0.037, +0.102]}
\newcommand{\expSizeCallsExploreVsComplexity}{+0.027 [$-$0.032, +0.086]}
\newcommand{\expSizeCallsConfirmVsSize}{+0.069 [+0.010, +0.135]}
\newcommand{\expSizeCallsExploreVsSize}{+0.028 [$-$0.021, +0.080]}
\newcommand{\expSizeCallsConfirmPrecision}{0.22}
\newcommand{\expSizeCallsConfirmRecall}{0.76}
\newcommand{\expSizeFanConfirmMcc}{0.280}
\newcommand{\expSizeFanExploreMcc}{0.317}
\newcommand{\expComplexityConfirmMcc}{0.254}
\newcommand{\expComplexityExploreMcc}{0.287}
\newcommand{\expSizeOnlyConfirmMcc}{0.217}
\newcommand{\expSizeOnlyExploreMcc}{0.286}
\newcommand{\expCallsAloneConfirmMcc}{0.214}
\newcommand{\expFanAloneConfirmMcc}{0.158}
\newcommand{\expDensityCallsConfirmMcc}{0.229}
\newcommand{\expDensityConfirmMcc}{0.039}
\newcommand{\expFormBSumConfirmMcc}{0.241}
\newcommand{\expRateConfirmMcc}{0.193}
\newcommand{\expXOConfirmVsSizeMax}{+0.036 [$-$0.031, +0.110]}
\newcommand{\expXOExploreVsSizeMax}{+0.016 [$-$0.050, +0.078]}
\newcommand{\expXOFontanaVsSizeMax}{+0.061 [+0.011, +0.110]}
\newcommand{\expFontanaTop}{fan-out (callee files) alone}
\newcommand{\expFontanaTopMcc}{0.847}
\newcommand{\expFontanaComplexityMcc}{0.835}
\newcommand{\expGateAnyConfirmMcc}{0.744}
\newcommand{\expGateAnyExploreMcc}{0.679}
\newcommand{\expGateAnyConfirmPrecision}{0.69}
\newcommand{\expGateAnyConfirmRecall}{0.87}
\newcommand{\expGateAnyConfirmVsTokens}{+0.119 [+0.070, +0.171]}
\newcommand{\expGateAnyExploreVsTokens}{+0.059 [$-$0.002, +0.118]}
\newcommand{\expGateAnyConfirmVsLines}{+0.224 [+0.164, +0.280]}
\newcommand{\expGateIntConfirmMcc}{0.697}
\newcommand{\expGateIntExploreMcc}{0.679}
\newcommand{\expGateIntConfirmVsTokens}{+0.072 [+0.009, +0.142]}
\newcommand{\expGateIntExploreVsTokens}{+0.059 [+0.004, +0.113]}
\newcommand{\expTokensConfirmMcc}{0.626}
\newcommand{\expTokensExploreMcc}{0.620}
\newcommand{\expLinesConfirmMcc}{0.520}
\newcommand{\expWeightConfirmMcc}{0.465}
\newcommand{\expWeightConfirmVsTokens}{$-$0.161 [$-$0.215, $-$0.110]}
\newcommand{\expDensityGateConfirmMcc}{0.017}
\newcommand{\expMethodDensityConfirmMcc}{$-$0.042}
\newcommand{\expGateLeaveConfirmMcc}{0.561}
\newcommand{\expToolPmdTop}{cyclomatic complexity, summed (WMC)}
\newcommand{\expToolPmdTopMcc}{0.582}
\newcommand{\expToolWmcTop}{cyclomatic complexity, summed (WMC)}
\newcommand{\expToolWmcTopMcc}{0.716}
\newcommand{\expToolDecorTop}{lines summed $\times$ number of methods (size only)}
\newcommand{\expToolDecorTopMcc}{0.455}
\newcommand{\expToolNpathCcMcc}{0.358}
\newcommand{\expToolNpathCcGateMcc}{0.339}
\newcommand{\expToolBrainCogMcc}{0.220}
\newcommand{\expToolBrainCogGateMcc}{0.210}
\newcommand{\matchedMadeUtc}{20261026 234905 UTC}
\newcommand{\matchedShaShort}{d6b77463f469}
\newcommand{\matchedClassFlags}{43{,}032}
\newcommand{\matchedClassUnits}{171{,}748}
\newcommand{\matchedClassRate}{25.1\%}
\newcommand{\matchedXAucConfirm}{0.768}
\newcommand{\matchedXAucExplore}{0.804}
\newcommand{\matchedOAucConfirm}{0.763}
\newcommand{\matchedOAucExplore}{0.774}
\newcommand{\matchedOAucVsXConfirm}{$-$0.005 [$-$0.043, +0.033]}
\newcommand{\matchedOAucVsXExplore}{$-$0.030 [$-$0.055, $-$0.006]}
\newcommand{\matchedNAucVsXConfirm}{+0.002 [$-$0.027, +0.035]}
\newcommand{\matchedDAucVsXConfirm}{$-$0.137 [$-$0.195, $-$0.078]}
\newcommand{\matchedSizeCallsAucVsXConfirm}{+0.008 [$-$0.028, +0.048]}
\newcommand{\matchedSizeCallsAucVsXExplore}{$-$0.011 [$-$0.038, +0.014]}
\newcommand{\matchedXOwnFlags}{117}
\newcommand{\matchedOOwnFlags}{301}
\newcommand{\matchedXOwnMcc}{0.187}
\newcommand{\matchedXMatchedMcc}{0.240}
\newcommand{\matchedXMatchedRecall}{0.64}
\newcommand{\matchedXMatchedPrecision}{0.21}
\newcommand{\matchedOMcc}{0.272}
\newcommand{\matchedOVsXConfirm}{+0.032 [$-$0.036, +0.099]}
\newcommand{\matchedOVsXExplore}{$-$0.015 [$-$0.067, +0.035]}
\newcommand{\matchedNVsXConfirm}{+0.052 [$-$0.005, +0.110]}
\newcommand{\matchedSizeCallsVsXConfirm}{+0.046 [$-$0.018, +0.111]}
\newcommand{\matchedSizeCallsVsXExplore}{$-$0.002 [$-$0.060, +0.055]}
\newcommand{\matchedDVsXConfirm}{$-$0.078 [$-$0.150, $-$0.003]}
\newcommand{\matchedFontanaOAucVsX}{+0.021 [+0.004, +0.043]}
\newcommand{\matchedFontanaOVsX}{+0.078 [+0.026, +0.134]}
\newcommand{\matchedFontanaSizeCallsAucVsX}{+0.023 [+0.006, +0.044]}
\newcommand{\matchedGateAucConfirm}{0.975}
\newcommand{\matchedGateAucExplore}{0.953}
\newcommand{\matchedGateAucVsXmConfirm}{+0.039 [+0.013, +0.063]}
\newcommand{\matchedGateAucVsXmExplore}{+0.059 [+0.024, +0.096]}
\newcommand{\matchedTokensAucConfirm}{0.963}
\newcommand{\matchedTokensAucVsXmConfirm}{+0.027 [$-$0.000, +0.054]}
\newcommand{\matchedTokensAucVsXmExplore}{+0.068 [+0.032, +0.104]}
\newcommand{\matchedCogAucConfirm}{0.936}
\newcommand{\matchedIntGateAucVsXmConfirm}{$-$0.081 [$-$0.145, $-$0.025]}
\newcommand{\matchedIntGateAucConfirm}{0.855}
\newcommand{\matchedWeightAucVsXmConfirm}{0.000 [$-$0.033, +0.029]}
\newcommand{\matchedGateVsXmConfirm}{+0.054 [$-$0.002, +0.105]}
\newcommand{\matchedGateVsXmExplore}{+0.064 [+0.011, +0.119]}
\newcommand{\matchedMethodRate}{14.1\%}
\newcommand{\matchedTokensVsXmConfirm}{+0.072 [+0.009, +0.133]}
\newcommand{\revMadeUtc}{20261026 235742 UTC}
\newcommand{\revShaShort}{321ac1632447}
\newcommand{\revLongReviews}{3{,}054}
\newcommand{\revBlobReviews}{3{,}518}
\newcommand{\revLongUnmapped}{303}
\newcommand{\revReviewers}{26}
\newcommand{\revNoCallLong}{3}
\newcommand{\revNoCallReviews}{300}
\newcommand{\revLeafTenRate}{0.4\%}
\newcommand{\revNonLeafTenRate}{1.0\%}
\newcommand{\revLeafTenN}{770}
\newcommand{\revNonLeafTenN}{823}
\newcommand{\revLeafTwentyFiveRate}{17.4\%}
\newcommand{\revNonLeafTwentyFiveRate}{24.3\%}
\newcommand{\revLeafTwentyFiveN}{230}
\newcommand{\revNonLeafTwentyFiveN}{571}
\newcommand{\revLeafFiftyRate}{80.0\%}
\newcommand{\revNonLeafFiftyRate}{81.6\%}
\newcommand{\revLeafFiftyN}{80}
\newcommand{\revNonLeafFiftyN}{365}
\newcommand{\revLeafHundredRate}{94.1\%}
\newcommand{\revNonLeafHundredRate}{94.0\%}
\newcommand{\revLeafHundredN}{17}
\newcommand{\revNonLeafHundredN}{151}
\newcommand{\revLeafOverRate}{85.7\%}
\newcommand{\revNonLeafOverRate}{95.0\%}
\newcommand{\revLeafOverN}{7}
\newcommand{\revNonLeafOverN}{40}
\newcommand{\revLongLeafMinusNonLeaf}{$-$0.033 [$-$0.141, +0.045]}
\newcommand{\revReviewersLeafLower}{4}
\newcommand{\revReviewersShown}{6}
\newcommand{\revDisagreeLongLeaf}{21.9\%}
\newcommand{\revDisagreeLongNonLeaf}{19.7\%}
\newcommand{\revDisagreeLongLeafN}{32}
\newcommand{\revDisagreeLongNonLeafN}{173}
\newcommand{\revCutoffReviewers}{10}
\newcommand{\revCutoffMin}{17}
\newcommand{\revCutoffMax}{34}
\newcommand{\revCutoffMedian}{21}
\newcommand{\revCutoffMccMin}{0.72}
\newcommand{\revCutoffMccMax}{0.91}
\newcommand{\revPosMedianMin}{30}
\newcommand{\revPosMedianMax}{44}
\newcommand{\revNegMedianMax}{15}
\newcommand{\revBlobHundredRate}{13.3\%}
\newcommand{\revBlobHundredN}{2{,}113}
\newcommand{\revBlobHundredDisagree}{38.0\%}
\newcommand{\revBlobTwoFiftyRate}{38.6\%}
\newcommand{\revBlobTwoFiftyN}{663}
\newcommand{\revBlobTwoFiftyDisagree}{61.7\%}
\newcommand{\revBlobFiveHundredRate}{48.9\%}
\newcommand{\revBlobFiveHundredN}{380}
\newcommand{\revBlobFiveHundredDisagree}{43.7\%}
\newcommand{\revBlobThousandRate}{43.1\%}
\newcommand{\revBlobThousandN}{246}
\newcommand{\revBlobThousandDisagree}{47.5\%}
\newcommand{\revBlobOverRate}{39.7\%}
\newcommand{\revBlobOverN}{116}
\newcommand{\revBlobOverDisagree}{41.7\%}
\newcommand{\revFlaggedLeafN}{100}
\newcommand{\revFlaggedRestN}{460}
\newcommand{\revFlaggedLeafMedianLen}{39}
\newcommand{\revFlaggedRestMedianLen}{42}
\newcommand{\revFlaggedLeafRate}{63.0\%}
\newcommand{\revFlaggedRestRate}{83.3\%}
\newcommand{\revBlobLargeCallsLowRate}{39.7\%}
\newcommand{\revBlobLargeCallsMidRate}{35.8\%}
\newcommand{\revBlobLargeCallsHighRate}{50.4\%}
\newcommand{\matchedGateAucVsTokensConfirm}{+0.011 [+0.002, +0.020]}
\newcommand{\matchedGatePctVsTokensConfirm}{+0.005 [$-$0.002, +0.013]}
\newcommand{\matchedGateMatchedVsTokensConfirm}{$-$0.018 [$-$0.056, +0.017]}
\newcommand{\matchedLinesPctVsTokensConfirm}{+0.005 [$-$0.002, +0.012]}
\newcommand{\matchedGatePctConfirm}{0.971}
\newcommand{\matchedLinesPctConfirm}{0.970}
\newcommand{\matchedTokensPctConfirm}{0.966}
\newcommand{\matchedCogPctConfirm}{0.934}
\newcommand{\matchedLinesOwnFlagsConfirm}{332}
\newcommand{\matchedTokensOwnFlagsConfirm}{252}
\newcommand{\matchedGateOwnFlagsConfirm}{157}
\newcommand{\matchedCogOwnFlagsConfirm}{110}
\newcommand{\matchedGateOwnMccConfirm}{0.744}
\newcommand{\matchedLinesOwnMccConfirm}{0.520}
\newcommand{\matchedGateAucVsTokensExplore}{$-$0.009 [$-$0.022, +0.003]}
\newcommand{\matchedGatePctVsTokensExplore}{$-$0.006 [$-$0.018, +0.003]}
\newcommand{\matchedGateMatchedVsTokensExplore}{$-$0.047 [$-$0.083, $-$0.010]}
\newcommand{\matchedLinesPctVsTokensExplore}{$-$0.002 [$-$0.011, +0.007]}
\newcommand{\matchedGatePctExplore}{0.954}
\newcommand{\matchedLinesPctExplore}{0.958}
\newcommand{\matchedTokensPctExplore}{0.960}
\newcommand{\matchedCogPctExplore}{0.892}
\newcommand{\matchedLinesOwnFlagsExplore}{358}
\newcommand{\matchedTokensOwnFlagsExplore}{265}
\newcommand{\matchedGateOwnFlagsExplore}{148}
\newcommand{\matchedCogOwnFlagsExplore}{124}
\newcommand{\matchedGateOwnMccExplore}{0.679}
\newcommand{\matchedLinesOwnMccExplore}{0.477}
\newcommand{\matchedOAucVsHvolConfirm}{$-$0.021 [$-$0.055, +0.014]}
\newcommand{\matchedSizeCallsAucVsHvolConfirm}{$-$0.008 [$-$0.032, +0.016]}
\newcommand{\logMadeUtc}{20261026 230649 UTC}
\newcommand{\logShaShort}{2fa130f412f6}
\newcommand{\logClassForms}{1813}
\newcommand{\logClassBetter}{64}
\newcommand{\logClassWorse}{308}
\newcommand{\logXFlagsRaw}{117}
\newcommand{\logXFlagsLog}{5}
\newcommand{\logXMccLog}{0.023}
\newcommand{\logOMccLog}{0.229}
\newcommand{\logOMinusRaw}{$-$0.043 [$-$0.094, $-$0.000]}
\newcommand{\logClassTopMcc}{0.295}
\newcommand{\logClassTopRaw}{0.273}
\newcommand{\logMethodForms}{280}
\newcommand{\logMethodBetter}{17}
\newcommand{\logMethodWorse}{100}
\newcommand{\logTokensFlagsRaw}{252}
\newcommand{\logTokensFlagsLog}{11}


\newcommand{\madz}{modified $z \geq 3.5$}
\newcommand{\fitwidth}[1]{\resizebox{\ifdim\width>\textwidth\textwidth\else\width\fi}{!}{#1}}
\newcommand{\sig}[1]{\texttt{#1}}
\newcommand{\Xs}{X$^*$}
\newcommand{\Xm}{X$^*_m$}
\newcommand{\Dflow}{D$_{\text{flow}}$}
\newcommand{\Oflow}{O$_{\text{flow}}$}
\newcommand{\Eflow}{E$_{\text{flow}}$}
\newcommand{\Aflow}{A$_{\text{flow}}$}
\newcommand{\dD}{$\Delta$D}
\newcommand{\dO}{$\Delta$O}
\newcommand{\dDO}{$\Delta$DO}

\title{\bfseries Does coupling divergence tell a god unit from a complex
algorithm?\\ A pre-registered test of weighting complexity concentration by
the data a unit hands on}

% Zenodo DOI, reserved before publishing and printed under the author block.
\newcommand{\doiline}{\ifdefined\thedoi\\\small\texttt{doi:\thedoi}\fi}

\author{Daniel Sagenschneider\\\small OfficeFloor\\\small\texttt{daniel@officefloor.net}\doiline}

\date{October 2026}

\begin{document}
\maketitle

\begin{abstract}
\noindent
A companion study found that a complexity outlier rule shows where
complexity concentrates but cannot tell a god unit from a complex
algorithm. This paper pre-registers a test of one way to tell them apart: weight a
unit's complexity by its \emph{coupling divergence}, the data it hands on
minus the data it receives, so that a leaf is set aside. Coupling is
measured by resolving every call in \mlcqMeasured{} MLCQ repositories,
\seedThreeMeasured{} fresh Java repositories and \fontanaSystems{} Qualitas
Corpus systems, and the weighted signal is scored by the same scale-free
outlier rule as complexity alone against MLCQ's reviewers, Fontana's God
Class and five tool definitions. The claim fails. The primary, the pooled
MCC of the weighted form minus complexity alone against human Blob, is
\primMicro{}, no difference resolved and underpowered, and every resolved
reading of the whole factor goes against it: it loses to the exemption
alone, to a weighting by the number of methods, and to complexity alone on
the threshold-free ranking. The one thing the data settle is at method
level. What MLCQ's reviewers call a Long Method is an absolute length: a
constant of \lwLengthCutoff{} lines or more, chosen on the explore half,
reaches MCC \lwLengthMccConfirm{} on the confirm half, \lwLengthVsRule{}
over the best per-repository form, each reviewer's own best cutoff lies
between \revCutoffMin{} and \revCutoffMax{} lines, and no form built from
the method's calls, variables, data flow or the cost of cutting it under
that length does better. The next form of the construct is not a detector
but a cost: what a long method's pieces must exchange to bring it under
the constant.
\end{abstract}

\section{Introduction}

A companion study \cite{paperone} tested a god-class and god-method
detector with no raw-scale constant: a unit is flagged when its complexity
is a robust outlier for its own codebase, a modified z-score of 3.5 or more
on a scale fitted to that codebase \cite{iglewicz1993}. The rule found where
the complexity concentrates, and its extra flags, read in the source, were
small classes doing the work they were written to do. A complexity measure
cannot separate a god unit from a complex algorithm because the difference
is not in the complexity: both hold more of it than their neighbours. The
difference, that study concluded, is in the coupling, and it named this one
as the test of that reading.

The construct tested here is \emph{coupling divergence}. For a method with
$p$ parameters that makes calls into its own repository, each passing
$a(c)$ arguments, the data coming in is what the method receives and the
data going out is what it hands on; the divergence is the difference. A
method that receives more than it sends is a \emph{catchment}, a leaf that
processes what it is given, which is where the complexity of an algorithm
belongs. A method that sends more than it receives is a \emph{junction},
and a unit that directs the data and also holds the logic is what people
call a god unit. The pre-registered claim is that multiplying a unit's
complexity by a factor that is zero for a catchment and grows with the
divergence of a junction recovers the units people call god classes and
methods better than complexity alone. The claim was committed, with its
estimand, decision rule, sample and bootstrap, before any coupling had been
measured on the confirmatory material.

This paper reports three things. The claim did not survive the test, and
the registered readings say where it fails. What MLCQ's reviewers call a
Long Method turned out to be an absolute length, \lwLengthCutoff{} lines or
more, which a per-repository rule is the wrong model for and which nothing
measured inside a method refines. And the construct has a next form, not as
a detector but as the cost of cutting a long method under that length,
which is stated as future work. The full registered report, every
secondary, sensitivity and table, is in the ancillary data; this paper
keeps what those three points need.

\section{Background}

The relative outlier rule, the signal families and the oracles are the
companion study's \cite{paperone}. A unit is a method (every function the
parser reports in a \texttt{.java} file) or a class (the file). A method's
lines are the code lines of its own body; the methods of an anonymous class
declared inside it are units of their own, so a short method that returns a
large anonymous class counts as short. Ten method signals in three families
(size: lines, tokens, Halstead volume; structure: cyclomatic
\cite{mcCabe1976}, cognitive \cite{campbell2018}, nesting, exit points;
density: cyclomatic density, maintainability index, parameters) give 31
class signals by aggregating each as a sum, a mean and a maximum over the
file's methods, plus the number of methods. The rule standardises a signal
within each repository by its median and median absolute deviation and
flags units at \madz{}; where more than half the units tie at the median
the rule falls back to the mean absolute deviation, which matters here
because the weighted forms are zero on every exempt unit.

Coupling as an ingredient of a god definition is old. Lanza and Marinescu's
God Class conjoins a weighted method count, low cohesion and access to
foreign data \cite{lanza2006, marinescu2004}, so agreement with PMD, which
implements it, is partly circular for any coupling measure. Henry and
Kafura's information flow multiplies a procedure's length by the square of
fan-in times fan-out \cite{henry1981}, and Kitchenham, Pickard and Linkman's
replication found size predicted as well \cite{kitchenham1990}, which is
why a size-only weighting is one of this paper's arms; El Emam and
colleagues showed that object-oriented metrics lose most of their predictive
power once class size is controlled \cite{elemam2001}. Divergence differs
from information flow in subtracting what comes in from what goes out.
Yourdon and Constantine's tramp data \cite{yourdon1979} is the ancestor of
the construct behind it, a \emph{sinew}: data that reaches the code that
needs it only by riding the call chain. In Tempero and Ralph's terms
\cite{tempero2026metrics} the construct is the sinew, the entity a method
or class, the attribute the direction of the data coupled to the control
flow, and divergence one indicator of it. Arif, Kuutila and Ralph's factor
analysis found in-coupling and out-coupling loading as separate constructs
\cite{arif2026}, which is why the outgoing side is read on its own as a
pre-registered secondary; Zhang, Wen and Tempero's typed dependency
distributions are the closest published typology to the in and out counted
here \cite{zhang2025deps}. Cognitive complexity counts a whole
\texttt{switch} once \cite{campbell2018}, the one place in the detection
literature where a long self-contained method is deliberately discounted.

The human oracle is MLCQ \cite{mlcq2020}, professional developers' reviews
of Java open-source classes and methods, with the companion study's
strict-majority aggregation and its explore and confirm halves frozen
before any repository was measured; classifiers trained on it reach a
median MCC of about 0.51 for Blob and 0.81 for Long Method \cite{mlcq2023},
a reference point rather than a baseline since the rules here see no label.
Fontana and colleagues' God Class instances on the Qualitas Corpus
\cite{fontana2016, tempero2010} are a second human oracle from a different
team on older code, read only as a paired difference because their
instances are balanced and extreme \cite{dinucci2018}.

\section{The claim and how it is read}
\label{sec:claim}

\paragraph{The measures.} For each method, under the primary \emph{flow}
definition of coupling, $\mathit{in} = 1 + p$ and $\mathit{out} = \sum_c
(1 + a(c))$ over the method's internal calls, those the symbol solver
resolves to a declaration in the repository's own sources; a call to the
JDK or to a dependency counts for nothing. Two secondary definitions,
\emph{data} (parameters and arguments only) and \emph{message} (the first
item as the message itself), run through every analysis. The divergence is
$d = \mathit{out} - \mathit{in}$. A class's $\mathit{out}$ sums its methods'
calls that leave the file and its $\mathit{in}$ the parameters of its
methods callable from outside the file.

\paragraph{The forms.} The divergence factor is $f_D = \max(d + 1, 0)$:
a catchment ($d < 0$) is 0, a pipe ($d = 0$) is 1, a junction with
divergence $j$ is $j + 1$. It is the product of the exemption $f_E = [d
\geq 0]$ and the amplification $f_A = \max(d + 1, 1)$, each measured
apart. The outgoing factor is $f_O = \mathit{out} + 1$, registered as a
secondary after an exploratory observation on agent-evolved code; the size
null $f_N$ is the class's number of methods, which carries no coupling.
Every class signal X gives X$\cdot$D, X$\cdot$E, X$\cdot$A, X$\cdot$O and
X$\cdot$N, each scored by the same rule on every unit of the repository,
exempt units as zeros, exactly as complexity alone.

\paragraph{The primary.} The base signal \Xs{} is chosen on the MLCQ explore
half by complexity alone, among the peak and density families, as the
signal whose flags maximise the pooled MCC against human Blob. The estimand
is the pooled MCC of \Xs$\cdot$\Dflow{} minus \Xs{}, both by \madz{},
against MLCQ human Blob over the reviewed classes of the confirm
repositories, as a paired difference with a 95\% percentile bootstrap over
repositories (\bootB{} resamples, seed \bootSeed{}). A lower bound above 0
reads \emph{refines}; an upper bound below 0, \emph{worse}; otherwise
\emph{no difference resolved}. An effective-sample gate, sized on the
explore half, marks the verdict \emph{underpowered} when fewer confirm
repositories carry a reviewed Blob than it requires. Beside it,
\Xs$\cdot$\Dflow{} is read against \Xs$\cdot$N, the size null.

\paragraph{The leaf diagnostic and the construct.} Among the units \Xs{}
flags, the human positive rate of catchments against the rest: the
exemption is \emph{supported} if catchments are called god units less often
(upper bound below 0), \emph{contradicted} if more often, otherwise
\emph{unresolved}. The consistency analysis then asks whether divergence
holds up across oracles and signal families: for every oracle, level and
family, three paired differences averaged over the family's signals, \dD{}
$= \mathrm{MCC}(\mathrm{X}\cdot\mathrm{D}) - \mathrm{MCC}(\mathrm{X})$,
\dO{} the same for X$\cdot$O, and \dDO{} between the two, one bootstrap
giving every interval. Four predictions were registered on flow: P1, \dD{}
$> 0$ for human Blob in the peak and density families; P2, \dD{} larger in
peak and structure than in size; P3, \dDO{} $\geq 0$ for human Blob,
against the competing expectation \dDO{} $\leq 0$ from the exploratory
observation; P4, \dD{} positive in a majority of the counted tool cells.
The verdict is \emph{supported} if P1 holds and P3 does not read
outgoing-only; \emph{outgoing data only} if P1 holds for \dO{} and P3 reads
outgoing-only; \emph{not supported} if neither \dD{} nor \dO{} resolves
above 0 for human Blob in any non-size family; otherwise \emph{unresolved}.

\paragraph{What was seen before the plan.} The companion study's main
result against PMD had been read; its human-label read-out had not. The
plan was amended ten times after registration, each amendment a commit
stating what had been seen; the first nine before any confirm or seed-3
repository had been measured, the tenth correcting a fault found after the
confirmatory results had been read (Section~\ref{sec:deviations}).

\section{Material and measurement}
\label{sec:design}

Three measured sets. (1) The \mlcqRepos{} MLCQ repositories at their
labelled commits, \mlcqMeasured{} measured; the confirm half holds
\mlcqConfirmMeasured{} repositories, of which \mlcqConfirmUsable{} measured
and pass the companion study's substance floor. MLCQ holds
\mlcqBlobReviewed{} reviewed classes with \mlcqBlob{} Blob positives and
\mlcqLongReviewed{} reviewed methods with \mlcqLong{} Long Method positives
after strict-majority aggregation, over both halves. (2) A fresh sample of
\seedThreeMeasured{} usable Java repositories drawn with seed 3 by the
companion study's sampler, disjoint from its earlier samples and from MLCQ,
carrying the tool oracles: PMD GodClass, an approximate DECOR Blob
\cite{moha2010, cktool}, WMC $\geq 47$, PMD NPath at method level and an
internal Brain Method. (3) Fontana's God Class: \fontanaInstances{}
manually evaluated instances on \fontanaSystems{} Qualitas Corpus system
versions, landing on \fontanaFilesReviewed{} reviewed files with
\fontanaFilesGod{} positives; the corpus's maintainer agreed to its use.
Test directories, build output and generated sources are dropped before any
rule is fitted.

Complexity is the companion study's extractor at its released commit;
coupling is a JavaParser measurer with its symbol solver, resolving against
the repository's sources and the JDK with no build and no dependencies,
under a time budget of eight hours per repository, which no repository
spent. The two are joined by file and start line (the \emph{joined} set,
which decides); every complexity field was also re-implemented on
JavaParser's tree so that both come from one parser (the
\emph{JavaParser-only} set), and every verdict is marked agrees or differs.
On the confirm half the measurer joined \joinMlcqShare{} of the extractor's
units; \callsMlcqInternal{} of calls resolved to the repository's own code,
\callsMlcqJdk{} to the JDK and \callsMlcqUnresolved{} not at all, almost
always a dependency's type, so $\mathit{out}$ reads low in proportion;
\diMlcqRepos{} of the usable confirm repositories declare a
dependency-injection framework, whose injected collaborators the parser
cannot see. Under flow, \roleMlcqClassCatchmentUnits{} of confirm classes
are catchments holding \roleMlcqClassCatchmentCc{} of the cyclomatic
complexity, and \roleMlcqMethodCatchmentUnits{} of methods, every method
with no internal call among them. Exempt units are zeros of the weighted
form and are more than half the units in the median repository
(\zeroMlcqDClass{} of classes), so the rule fell to its fallback scale in
\scaleMlcqDMeanAd{} of the \scaleMlcqClassRepos{} confirm repositories
against \scaleMlcqXMad{} for \Xs{} alone; the two arms of the primary are
therefore read under different scale estimators, and the arm fitted on the
non-exempt units only is a reported sensitivity.

\section{Results}

\subsection{The claim fails}
\label{sec:primary}

\begin{table}[t]
\centering
\caption{The primary read-out against MLCQ human Blob on the confirm half:
\primUniv{} reviewed classes, \primPos{} of them Blob, in \primRepos{}
repositories. Every arm is scored by \madz{} on every class of its
repository. Flagged and hits are pooled counts over the reviewed classes;
the paired difference is the pooled MCC of the arm minus that of \Xs{}, with
its 95\% percentile bootstrap interval over repositories.}
\label{tab:primary}
\small
\setlength{\tabcolsep}{4pt}
\fitwidth{% generated by paper/scripts/tables.py; do not edit
\begin{tabular}{lrrrrrrl}
\toprule
Arm (class form A, flow, non-private inputs) & MCC & Precision & Recall & FPR & Flagged & Hits & Paired difference against X$^*$ (95\% CI) \\
\midrule
X$^*$ = \texttt{nloc\_max}, complexity alone & 0.187 & 0.26 & 0.31 & 0.104 & 117 & 30 &  \\
X$^*\cdot$D$_{\text{flow}}$, the primary arm & 0.146 & 0.19 & 0.41 & 0.208 & 214 & 40 & $-$0.041 [$-$0.130, +0.054] \\
X$^*\cdot$E$_{\text{flow}}$, exemption only & 0.212 & 0.26 & 0.37 & 0.122 & 138 & 36 & +0.025 [$-$0.052, +0.100] \\
X$^*\cdot$A$_{\text{flow}}$, amplification only & 0.196 & 0.20 & 0.56 & 0.268 & 279 & 55 & +0.009 [$-$0.075, +0.104] \\
X$^*\cdot$N, the size null & 0.236 & 0.22 & 0.60 & 0.252 & 270 & 59 & +0.049 [$-$0.023, +0.121] \\
X$^*\cdot$D$_{\text{flow}}$ with constructors as inputs & 0.125 & 0.20 & 0.27 & 0.124 & 130 & 26 & $-$0.062 [$-$0.130, +0.011] \\
X$^*\cdot$O$_{\text{flow}}$, outgoing factor only & 0.272 & 0.23 & 0.69 & 0.279 & 301 & 68 & +0.085 [+0.001, +0.182] \\
\bottomrule
\end{tabular}
}
\end{table}

On the \xClassRepos{} explore repositories carrying a reviewed class,
complexity alone chose \sig{\xClass} as \Xs{} at a pooled MCC of
\xClassMcc{}; the explore half's own estimate set the gate at
\gateRequired{} confirm repositories with a reviewed Blob, and the confirm
half carries \primGateReposWithBlob{}, so the primary is underpowered by
the plan's rule whatever it reads. Table~\ref{tab:primary} gives every arm.
Complexity alone flags \primXFlagged{} of the \primUniv{} reviewed confirm
classes and hits \primXTp{} of the \primPos{} Blobs: MCC \primXMcc{}. The
primary arm flags \primDFlagged{} and hits \primDTp{}: MCC \primDMcc{}. The
paired difference is \textbf{\primMicro{}}, verdict \textbf{\primVerdict}.
The point estimate is on the wrong side and the hypothesised shape is
inverted: the factor gains recall (\primDRecall{} against \primXRecall{})
and loses precision (\primDPrecision{} against \primXPrecision{}).
Threshold-free, the modified z of \Xs{} ranks Blobs above non-Blobs at an
AUC of \primAucX{} and the weighted form at \primAucD{}, \primAucDiff{},
complexity alone the better ranking. The halves say which part costs: the
exemption alone is \primEMicro{} against \Xs{} and the amplification alone
\primAmpMicro{}, neither resolved, but the whole factor is resolved worse
than the exemption alone, \primDvsEMicro{}. Against the size null, a
weighting by the number of methods that carries no coupling, the coupling
factor is \textbf{\primDvsNMicro{}}: the rival explanation the plan
recorded, that a Blob is a big class with one complex method, beats the
construct by a resolved margin. The JavaParser-only set agrees on every
verdict (primary \primJpMicro{}).

\begin{table}[t]
\centering
\caption{The leaf diagnostic: among the reviewed units \Xs{} (or \Xm{})
flags, how many are catchments ($d < 0$), how many of those the reviewers
call positive, the same for the rest, and the paired difference in positive
rate (catchments minus the rest) with its interval over repositories.}
\label{tab:leaf}
\small
\setlength{\tabcolsep}{4pt}
\fitwidth{% generated by paper/scripts/tables.py; do not edit
\begin{tabular}{llrrrrll}
\toprule
Oracle & Definition, inputs & Catchments flagged & of them positive & Rest flagged & of them positive & Difference in positive rate (95\% CI) & Reading \\
\midrule
MLCQ Blob & flow, non-private & 38 & 8 & 79 & 22 & $-$0.068 [$-$0.238, +0.102] & unresolved \\
MLCQ Blob & flow, public & 33 & 7 & 84 & 23 & $-$0.062 [$-$0.229, +0.118] & unresolved \\
MLCQ Blob & flow, non-private, with constructors & 41 & 8 & 76 & 22 & $-$0.094 [$-$0.260, +0.069] & unresolved \\
MLCQ Blob & flow, public, with constructors & 34 & 7 & 83 & 23 & $-$0.071 [$-$0.236, +0.105] & unresolved \\
MLCQ Blob & data, non-private & 44 & 8 & 73 & 22 & $-$0.120 [$-$0.284, +0.041] & unresolved \\
MLCQ Blob & data, public & 34 & 8 & 83 & 22 & $-$0.030 [$-$0.196, +0.145] & unresolved \\
MLCQ Blob & data, non-private, with constructors & 52 & 10 & 65 & 20 & $-$0.115 [$-$0.280, +0.047] & unresolved \\
MLCQ Blob & data, public, with constructors & 44 & 10 & 73 & 20 & $-$0.047 [$-$0.213, +0.122] & unresolved \\
MLCQ Blob & message, non-private & 37 & 7 & 80 & 23 & $-$0.098 [$-$0.266, +0.060] & unresolved \\
MLCQ Blob & message, public & 31 & 5 & 86 & 25 & $-$0.129 [$-$0.296, +0.031] & unresolved \\
MLCQ Blob & message, non-private, with constructors & 38 & 7 & 79 & 23 & $-$0.107 [$-$0.278, +0.049] & unresolved \\
MLCQ Blob & message, public, with constructors & 33 & 6 & 84 & 24 & $-$0.104 [$-$0.269, +0.054] & unresolved \\
MLCQ Long Method & flow, all calls & 28 & 13 & 82 & 63 & $-$0.304 [$-$0.541, $-$0.082] & supported (catchments are called positive less often) \\
MLCQ Long Method & data, all calls & 28 & 15 & 82 & 61 & $-$0.208 [$-$0.413, $-$0.017] & supported (catchments are called positive less often) \\
MLCQ Long Method & message, all calls & 27 & 12 & 83 & 64 & $-$0.327 [$-$0.557, $-$0.102] & supported (catchments are called positive less often) \\
Fontana God Class & flow, non-private & 20 & 15 & 98 & 93 & $-$0.199 [$-$0.362, $-$0.001] & supported (catchments are called positive less often) \\
Fontana God Class & data, non-private & 19 & 15 & 99 & 93 & $-$0.150 [$-$0.316, +0.050] & unresolved \\
Fontana God Class & message, non-private & 22 & 17 & 96 & 91 & $-$0.175 [$-$0.322, 0.000] & unresolved \\
\bottomrule
\end{tabular}
}
\end{table}

The exemption, the claim in the question, has its own test
(Table~\ref{tab:leaf}). Against MLCQ Blob it is \leafReading{}:
\leafCatchRate{} of the flagged catchments are Blobs against
\leafRestRate{} of the rest, \leafDiff{}, every point estimate on the
exemption's side. Against MLCQ Long Method it is \leafMethodReading{}
(\leafMethodCatchRate{} against \leafMethodRestRate{}, \leafMethodDiff{})
and against Fontana's God Class \leafFontanaReading{} (\leafFontanaDiff{}).
A complex unit that receives more than it sends is called a god unit less
often by two of the three human oracles. The distinction is in the labels;
the factor built on it does not improve detection, because what the factor
adds to the exemption, an amplification by the divergence of a junction,
costs more than the exemption gains.

\begin{table}[t]
\centering
\caption{The consistency cells of the human oracles under the flow
definition, every signal of the family averaged: \dD{} is the paired MCC
difference of X$\cdot$D against X, \dO{} of X$\cdot$O against X, and \dDO{}
of X$\cdot$D against X$\cdot$O. One bootstrap over repositories gives every
interval; $n$ is the number of repositories (systems) carrying the oracle.}
\label{tab:cells-human}
\small
\setlength{\tabcolsep}{4pt}
\fitwidth{% generated by paper/scripts/tables.py; do not edit
\begin{tabular}{llllrlll}
\toprule
Oracle & Family & Relation & Counted & $n$ & $\Delta$D (95\% CI) & $\Delta$O (95\% CI) & $\Delta$DO (95\% CI) \\
\midrule
MLCQ Blob (class) & size & non-ingredient & yes & 208 & $-$0.032 [$-$0.114, +0.048] & +0.048 [$-$0.014, +0.108] & $-$0.080 [$-$0.144, $-$0.014] \\
MLCQ Blob (class) & peak & non-ingredient & yes & 208 & +0.041 [$-$0.031, +0.114] & +0.125 [+0.054, +0.195] & $-$0.084 [$-$0.139, $-$0.029] \\
MLCQ Blob (class) & density & non-ingredient & yes & 208 & +0.094 [+0.013, +0.174] & +0.151 [+0.076, +0.227] & $-$0.058 [$-$0.110, $-$0.009] \\
MLCQ Long Method (method) & size & non-ingredient & yes & 202 & $-$0.232 [$-$0.317, $-$0.144] & $-$0.159 [$-$0.200, $-$0.121] & $-$0.073 [$-$0.156, +0.012] \\
MLCQ Long Method (method) & structure & non-ingredient & yes & 202 & $-$0.044 [$-$0.115, +0.028] & $-$0.028 [$-$0.086, +0.034] & $-$0.017 [$-$0.061, +0.024] \\
MLCQ Long Method (method) & density & non-ingredient & yes & 202 & +0.280 [+0.197, +0.362] & +0.222 [+0.162, +0.280] & +0.058 [+0.006, +0.110] \\
Fontana God Class (class) & size & non-ingredient & yes & 72 & $-$0.073 [$-$0.193, +0.037] & +0.063 [$-$0.004, +0.121] & $-$0.136 [$-$0.243, $-$0.038] \\
Fontana God Class (class) & peak & non-ingredient & yes & 72 & +0.169 [+0.068, +0.269] & +0.315 [+0.249, +0.381] & $-$0.146 [$-$0.230, $-$0.067] \\
Fontana God Class (class) & density & non-ingredient & yes & 72 & +0.413 [+0.326, +0.503] & +0.573 [+0.499, +0.658] & $-$0.160 [$-$0.247, $-$0.079] \\
\bottomrule
\end{tabular}
}
\end{table}

The construct's verdict, by the pre-registered rule, is
\textbf{\consVerdict} (Table~\ref{tab:cells-human}): P1 \pOneHolds{} for
\dD{} (\pOneDPeak{} in peak, \pOneDDensity{} in density) and holds for
\dO{} (\pOneOPeak{}, \pOneODensity{}); P2 \pTwoHolds{}, but because the
factor damages the size families outright (\cellBlobSizeD{} for Blob); P3
reads \pThreeReads{}, every family's upper bound below 0, so subtracting
what comes in costs agreement; P4, \pFourPositive{} of \pFourCells{} counted
tool cells positive. The verdict is a statement about the operating point.
Read threshold-free and at a matched flag rate, post hoc, the outgoing
factor ranks Blobs no better than complexity alone on MLCQ (AUC
\matchedOAucConfirm{} against \matchedXAucConfirm{}) and gains by flagging
two and a half times as many classes; on Fontana alone it does rank better.
The secondaries, sensitivities, tool oracles and the Fontana replication,
all in the ancillary data, change nothing: the whole factor makes every
tool definition worse by a resolved margin at class level, loses at method
level against Long Method (\methMicro{}) through the amplification, moves
further against the construct when the dependency-injection repositories
are removed (\sensDiMicro{}), and returns to complexity alone, not past it,
when fitted on the non-exempt units only (\sensNeVsXMicro{}).

\subsection{What the reviewers call a Long Method is a length}
\label{sec:explore-forms}

Everything in this subsection was computed after every oracle of the study
had been read, by scripts that are in the ancillary data with their outputs
and digests; it is exploration and changes nothing above. The method-level
reading of the registered analysis was that length ranks Long Methods best
of every form, better than the registered cognitive signal, and that no form
built from calls ranks better than length does: the gate on any call that
looked like the exemption's own form turns out to zero the short trivial
methods and so recalibrate the per-repository constant, ranking no better
than length itself (percentile AUC against tokens
\matchedGatePctVsTokensConfirm{} on confirm).

\begin{table}[t]
\centering
\caption{Exploratory, post hoc: MLCQ's individual Long Method reviews on
measured units (\revLongReviews{} reviews by \revReviewers{} reviewers,
both halves), by the method's length and whether it is a leaf (no internal
call): the share of reviews that call the method a Long Method at any
severity.}
\label{tab:reviewers}
\small
\setlength{\tabcolsep}{4pt}
\fitwidth{% generated by paper/scripts/tables.py; do not edit
\begin{tabular}{lrrrr}
\toprule
Lines of the method & Leaf reviews & called Long Method & Non-leaf reviews & called Long Method \\
\midrule
1-10 & 770 & 0.4\% & 823 & 1.0\% \\
11-25 & 230 & 17.4\% & 571 & 24.3\% \\
26-50 & 80 & 80.0\% & 365 & 81.6\% \\
51-100 & 17 & 94.1\% & 151 & 94.0\% \\
101+ & 7 & 85.7\% & 40 & 95.0\% \\
\bottomrule
\end{tabular}
}
\end{table}

\paragraph{What the reviewers judge.} MLCQ keeps every review
(Table~\ref{tab:reviewers}). The share of reviews calling a method a Long
Method is \revNonLeafTenRate{} up to 10 lines,
\revNonLeafTwentyFiveRate{} at 11 to 25, \revNonLeafFiftyRate{} at 26 to 50
and \revNonLeafHundredRate{} at 51 to 100: a step at about twenty-five
lines. Each of the \revCutoffReviewers{} reviewers with ten or more
positives has a best length cutoff between \revCutoffMin{} and
\revCutoffMax{} lines (median \revCutoffMedian{}), an absolute constant
carried from reviewer to reviewer within a few lines, which is why a cutoff
tuned on MLCQ beat the per-repository rule in the companion study and why
no per-repository form here can reach it. At equal length a leaf is called
a Long Method no less often than a non-leaf (\revLeafFiftyRate{} against
\revNonLeafFiftyRate{} at 26 to 50 lines; over all long methods
\revLongLeafMinusNonLeaf{}). The exemption the leaf diagnostic supported is
narrower than the construct: among the reviewed methods cognitive
complexity flags, leaves are called Long Method less often than the rest
(\revFlaggedLeafRate{} against \revFlaggedRestRate{}) at the same median
length, so it is the branchy leaf that reviewers decline to call a Long
Method, not a leaf as such.

\paragraph{The constant, read once.} The constant the reviews imply was
read as the plan reads a cutoff: chosen on the explore half among absolute
cutoffs on lines and read on the confirm half (\lwConfirmRepos{}
repositories, \lwConfirmUnits{} reviewed methods, \lwConfirmPositives{}
positives). The cutoff chosen is \lwLengthCutoff{} lines or more. On confirm
it flags \lwLengthFlagged{} methods with \lwLengthHits{} hits, MCC
\lwLengthMccConfirm{} (precision \lwLengthPrecision{}, recall
\lwLengthRecall{}), \lwLengthVsRule{} over the best per-repository form,
the registered rule on length gated on any internal call. The MCC is flat
from 20 to 30 lines (\lwNlocTwentyMcc{} at 20, \lwNlocThirtyMcc{} at 30)
and has fallen by 40 (\lwNlocFortyMcc{}), so the constant is a mid-twenties
one rather than a sharp number.

\paragraph{What was tried against it.} Two further post hoc readings were
made on the same half, each labelled as such, to ask whether anything a
method's body contains refines that constant. The first scored
\lwHeldOutForms{} forms built from the method's variables, the live width a
reader holds at a line in several definitions, the counts of its locals,
parameters and fields, and the flow of data from what enters the method to
the calls it makes, by the same procedure; every one is resolved below
length (\lwFormsBelowLength{} of \lwHeldOutForms{}), the best of them at
\lwBestFormMcc{}, \lwBestFormVsLength{}. The second scored the cost of
cutting a long method into pieces under the constant, the parameters and
return values the pieces would exchange; its cheapest-seam form reaches
\splitCostMcc{}, \splitCostVsLength{} against the constant, and among the
\splitWithinN{} long methods, the only ground such a form can gain, it
ranks the \splitWithinPositives{} positives worse than length does (AUC
\splitCostAucWithin{} against \splitNlocAucWithin{}, \splitCostAucVsNloc{}).
The share of a class's pieces bound to the class, read as a god-class form,
is last on every class oracle (\splitLockedShareMcc{} against
\splitBlobLengthMcc{} for the longest method on the confirm half;
\splitFontanaLockedMcc{} against \splitFontanaHvolMcc{} for Halstead volume
on Fontana). Nothing in a long method's body, as these measures see it,
says which of the methods length flags the reviewers passed.

\section{Deviations from the plan}
\label{sec:deviations}

The full account, with the numbers of every analysis set aside, is in the
ancillary data. The human populations were first scored on the repositories
with a positive only: the harness attached a human oracle to a repository
only when it held a positive-labelled unit, so a repository whose reviewed
units are all negatives never entered the population. The fault was found
after a confirmatory analysis had been written, committed, read and
reported to the author in a draft of this paper; that is the one departure
from the plan's order, recorded as Amendment 10 with what had been seen.
The fix has no free parameter, and the explore-half choices and the
confirmatory analysis were then made again by the registered procedure,
the earlier files set aside with their digests; it restored
\restoredBlobRepos{} confirm repositories and \restoredBlobUnits{} reviewed
classes to the Blob population, and on the set-aside analysis the primary
read \priorPrimMicro{}, \priorPrimVerdict{}, with the construct's verdict
\priorConsVerdict{}. The confirmatory analysis was written three times in
all, the first set aside by Amendment 9 before any of its results was read.
Nine amendments preceded any confirm measurement, each committed with a
statement of what had been seen. The gate was not met: the confirm half
carries \primGateReposWithBlob{} repositories with a reviewed Blob against
the \gateRequired{} required, and MLCQ has no more, so a difference of 0.10
could have been missed on the primary.

\section{Discussion}
\label{sec:disc}

\paragraph{Where the construct fails.} Three things were expected of the
divergence factor: that it would raise precision by setting complex leaves
aside, that the leaves set aside would be called god units less often, and
that what a unit receives would matter as well as what it sends. The second
is seen, in two human oracles of three, and the reviews narrow it to the
branchy leaves complexity flags. The first and the third are not: precision
falls, recall rises, and in every family of every class-level oracle,
subtracting the inputs from the outputs costs agreement. A factor of $d +
1$ is unbounded and, for a class, grows with the number of calls the class
makes, so the modified z of X$\cdot$D is dominated by the classes that call
most whatever their complexity; and a class with a wide callable surface is
one that other code uses, so subtracting its surface pushes the hubs people
call Blobs towards the catchment side, the opposite of the construct's
reading. In Tempero and Ralph's terms \cite{tempero2026metrics}, divergence
as an indicator of the sinew construct fails convergent validity against
three human definitions of a god unit, and the component that converges at
the operating point is out-coupling, which Arif, Kuutila and Ralph found
loading as a construct of its own \cite{arif2026}; read threshold-free, even
that is a flag rate on MLCQ. The human Blob is weakly predicted by any of
this (complexity alone \primXMcc{}, classifiers trained on the labels 0.51
\cite{mlcq2023}), and part of the ceiling is the oracle.

\paragraph{What a Long Method is.} At method level the reviews settle more
than the forms do. What MLCQ's reviewers label a Long Method is a method of
\lwLengthCutoff{} lines or more, leaf or not, each reviewer within a few
lines of the others whatever the repository; the right detector for that
label is a constant on length, which the companion study's tuned cutoff
already reaches at 0.785 and which reaches \lwLengthMccConfirm{} here, and
the per-repository rule is the wrong model for it. Whatever else the term
god method means, the registered construct did not capture it, and nor did
the live width, the variable counts, the data flow or the split cost read
after it. What remains to register at method level is whether the
branchy-leaf exemption the reviews do show adds anything to a constant on
length, and whether the per-repository rule has a place on code whose scale
has moved, which is the purpose the plan stated and this data cannot test.
At class level the exploration leaves one post hoc hypothesis, total lines
times outgoing calls against its size-only counterpart, about +0.03 over
complexity alone on each half and unresolved; it is stated in the ancillary
data and is not a result of this study.

\paragraph{The next form of the sinew.} The construct's next form at method
level is not a detector but a cost. With the trigger fixed at the constant,
the question is what it takes to bring a long method under it: cut the
method into pieces and count what crosses each seam, the parameters a piece
must be handed and the value it must return, with a second return value
costing a record type. That is the sinew of the refactoring rather than of
the method, and it is the measure a tool would use, since length is the
trigger and the split is the proposal. Prior tools rank extract-method
candidates within a method \cite{tsantalis2011, haas2016}; none takes the
length cap as a constraint with a cost on what the pieces exchange, and
that instrument is to be built and read descriptively, on human-written and
on agent-written code, before any claim is registered for it.

\section{Threats to validity}
\label{sec:threats}

\paragraph{Construct.} Divergence counts parameters and arguments, not what
travels in them; return values and exceptions are not counted; a class's
inputs count the parameters of its callable methods, so a wide interface
reads as receiving much, which the discussion identifies as the likely
failure of the sign. A different definition of a class's inputs is a
different construct and was not registered.

\paragraph{Internal.} The populations are the registered ones only because
a fault was fixed and the analysis made again after its first results had
been read; the fix had no free parameter, and every headline reading before
and after it is in the ancillary data. The two arms of the primary are
fitted under different scale estimators. Unresolved calls make
$\mathit{out}$ a lower bound, and injected collaborators are invisible to
static parsing.

\paragraph{External.} Java only. MLCQ's repositories are concentrated in a
few organisations (\mlcqOwnerTopShare{} of the confirm half in ten owners)
and Fontana's instances are extreme. The purpose stated in the plan, a
definition that moves with the codebase as agent-maintained code changes
what a normal level of complexity is, was not tested: no agent-written
corpus with human labels exists at the scale a pre-registered test needs.

\paragraph{Statistical.} The primary is underpowered by the plan's own
gate, and no secondary is corrected for multiplicity. The intervals are
percentile bootstraps over repositories, coarse at the tails. The post hoc
readings of Section~\ref{sec:explore-forms} were made on the confirm half
after the constant on length had been chosen on explore and read on it, so
that half is a hold-out for the constant alone and not for anything chosen
after it.

\section{Reproducibility and data availability}
\label{sec:repro}

Every number in this paper is produced by a script from the harness's
analysis files, and every table is a generated fragment. The ancillary data
holds the full registered report (every secondary, sensitivity and table
this paper leaves out), the confirmatory analysis and the two that were set
aside with their digests as committed, the full account of every deviation
from the plan, the explore-half decisions, the measurer's validation, the
per-set run records, the sample files that pin every repository and system
to its commit or archive, the converted labels, and the post hoc
exploration's scripts, outputs and digests, labelled as such. The ancillary
README names the public repositories of the harness, the companion study,
the sampler and the measurer, the commit of each, and the command behind
every number.

\section{Conclusion}

Weighting a unit's complexity by its coupling divergence does not recover
the classes and methods people call god units better than complexity
alone. Against MLCQ's reviewers the primary resolves no difference and is
underpowered, and every resolved reading of the whole factor goes the other
way: it loses to the exemption alone, to a weighting by the number of
methods, and to complexity alone on the threshold-free ranking. The
exemption the claim rests on is partly in the labels, and what the factor
adds to it costs more than it gains. What the data settle is at method
level: MLCQ's reviewers call a method a Long Method at \lwLengthCutoff{}
lines or more, each within a few lines of the others, leaf or not; a
constant on length, MCC \lwLengthMccConfirm{} on the confirm half, is the
detector that matches them, and nothing read inside a method, its calls,
its live width, its variables, its data flow or the cost of cutting it,
refines that constant. The form of sinew left to build is the cost of the
cut.

\section*{Competing interests}

The author created and maintains OfficeFloor, a framework whose design
rests on taking data off the control flow, and the sinew construct that
divergence was meant to indicate comes from that line of work; a positive
result would have supported the framework's premise. The result is
negative on the registered claim, and the reading that survives, that
outgoing data carries the signal, was recorded as a competing expectation
before the data existed, from the author's own exploratory runs. The
fault in the harness's populations (Section~\ref{sec:deviations}) was
found after the confirmatory results had been read; the fix is the one the
plan's text dictates, and every headline reading before and after it is
reported; the plan, the seed, the decision rule
and the gate were fixed before the data existed, and a replication needs no
cooperation from the author. The companion study, the sampler and the
measurer are the author's.

\section*{Funding}

This work received no external funding and was carried out and paid for by
the author.

\begin{thebibliography}{99}

\bibitem{arif2026} M.~Arif, M.~Kuutila and P.~Ralph,
  ``A factor analysis of class-level software metrics,''
  in \emph{Proc. ACM/IEEE International Symposium on Empirical Software
  Engineering and Measurement (ESEM)}, 2026; arXiv:2609.20411.

\bibitem{fontana2016} F.~Arcelli Fontana, M.~V. M\"antyl\"a, M.~Zanoni and
  A.~Marino,
  ``Comparing and experimenting machine learning techniques for code smell
  detection,''
  \emph{Empirical Software Engineering}, vol.~21, no.~3, pp.~1143--1191, 2016.

\bibitem{cktool} M.~Aniche,
  ``Java code metrics calculator (CK),'' 2015. Version 0.7.0.

\bibitem{ck1994} S.~R. Chidamber and C.~F. Kemerer,
  ``A Metrics Suite for Object Oriented Design,''
  \emph{IEEE Transactions on Software Engineering}, vol.~20, no.~6,
  pp.~476--493, 1994.

\bibitem{campbell2018} G.~A. Campbell,
  ``Cognitive Complexity: a new way of measuring understandability,''
  SonarSource white paper, 2018.

\bibitem{dinucci2018} D.~Di~Nucci, F.~Palomba, D.~A. Tamburri, A.~Serebrenik
  and A.~De~Lucia,
  ``Detecting code smells using machine learning techniques: Are we there
  yet?''
  in \emph{Proc. 25th IEEE International Conference on Software Analysis,
  Evolution and Reengineering (SANER)}, pp.~612--621, 2018.

\bibitem{elemam2001} K.~El~Emam, S.~Benlarbi, N.~Goel and S.~N. Rai,
  ``The confounding effect of class size on the validity of object-oriented
  metrics,''
  \emph{IEEE Transactions on Software Engineering}, vol.~27, no.~7,
  pp.~630--650, 2001.

\bibitem{friedman1937} M.~Friedman,
  ``The use of ranks to avoid the assumption of normality implicit in the
  analysis of variance,''
  \emph{Journal of the American Statistical Association}, vol.~32, no.~200,
  pp.~675--701, 1937.

\bibitem{haas2016} R.~Haas and B.~Hummel,
  ``Deriving extract method refactoring suggestions for long methods,'' in
  \emph{Software Quality Days (SWQD 2016)}, LNBIP vol.~238, Springer, 2016,
  pp.~144--155.

\bibitem{henry1981} S.~Henry and D.~Kafura,
  ``Software structure metrics based on information flow,''
  \emph{IEEE Transactions on Software Engineering}, vol.~SE-7, no.~5,
  pp.~510--518, 1981.

\bibitem{iglewicz1993} B.~Iglewicz and D.~C. Hoaglin,
  \emph{How to Detect and Handle Outliers}. ASQC Quality Press, 1993.

\bibitem{kitchenham1990} B.~A. Kitchenham, L.~M. Pickard and S.~J. Linkman,
  ``An evaluation of some design metrics,''
  \emph{Software Engineering Journal}, vol.~5, no.~1, pp.~50--58, 1990.

\bibitem{lanza2006} M.~Lanza and R.~Marinescu,
  \emph{Object-Oriented Metrics in Practice}. Springer, 2006.

\bibitem{mlcq2020} L.~Madeyski and T.~Lewowski,
  ``MLCQ: Industry-relevant code smell data set,''
  in \emph{Proc. 24th International Conference on Evaluation and Assessment in
  Software Engineering (EASE)}, pp.~342--347, 2020. Data set version 1.1,
  Zenodo record 3666840.

\bibitem{mlcq2023} L.~Madeyski and T.~Lewowski,
  ``Detecting code smells using industry-relevant data,''
  \emph{Information and Software Technology}, vol.~155, 107112, 2023.

\bibitem{marinescu2004} R.~Marinescu,
  ``Detection strategies: metrics-based rules for detecting design flaws,''
  in \emph{Proc. 20th IEEE International Conference on Software Maintenance
  (ICSM)}, pp.~350--359, 2004.

\bibitem{mcCabe1976} T.~J. McCabe,
  ``A Complexity Measure,''
  \emph{IEEE Transactions on Software Engineering}, vol.~SE-2, no.~4,
  pp.~308--320, 1976.

\bibitem{moha2010} N.~Moha, Y.-G. Gu\'eh\'eneuc, L.~Duchien and A.-F. Le~Meur,
  ``DECOR: A Method for the Specification and Detection of Code and Design
  Smells,''
  \emph{IEEE Transactions on Software Engineering}, vol.~36, no.~1, pp.~20--36,
  2010.

\bibitem{paperone} D.~Sagenschneider,
  ``Complexity concentration as a god-class and god-method detector: a
  pre-registered test of a scale-free outlier rule against tool definitions
  and human labels,'' 2026. \url{https://doi.org/10.5281/zenodo.23581333}

\bibitem{schumacher2010} J.~Schumacher, N.~Zazworka, F.~Shull, C.~Seaman and
  M.~Shaw,
  ``Building empirical support for automated code smell detection,''
  in \emph{Proc. ACM-IEEE International Symposium on Empirical Software
  Engineering and Measurement (ESEM)}, article 8, 2010.

\bibitem{tempero2010} E.~Tempero, C.~Anslow, J.~Dietrich, T.~Han, J.~Li,
  M.~Lumpe, H.~Melton and J.~Noble,
  ``The Qualitas Corpus: A curated collection of Java code for empirical
  studies,''
  in \emph{Proc. 17th Asia Pacific Software Engineering Conference (APSEC)},
  pp.~336--345, 2010.

\bibitem{tempero2026metrics} E.~Tempero and P.~Ralph,
  ``Making software metrics useful,''
  \emph{IEEE Computer}, in press, 2026; arXiv:2603.16012.

\bibitem{tsantalis2011} N.~Tsantalis and A.~Chatzigeorgiou,
  ``Identification of extract method refactoring opportunities for the
  decomposition of methods,'' \emph{Journal of Systems and Software},
  vol.~84, no.~10, pp.~1757--1782, 2011.

\bibitem{yourdon1979} E.~Yourdon and L.~L. Constantine,
  \emph{Structured Design: Fundamentals of a Discipline of Computer Program
  and Systems Design}. Prentice-Hall, 1979.

\bibitem{zhang2025deps} Z.~Zhang, E.~Wen and E.~Tempero,
  ``Analyzing variations in dependency distributions due to code smell
  interactions,'' arXiv:2509.03896, 2025.

\end{thebibliography}

\end{document}
